\BourbakiExerciseGroup

\begin{enumerate}
\def\labelenumi{\arabic{enumi})}
\item
  Let \(\mathcal{C}_1, \mathcal{C}_2\) be two topologies on the same set X. Show that \(\mathcal{C}_1\) is finer than \(\mathcal{C}_2\) if and only if every filter on X which is convergent in the topology \(\mathcal{C}_1\) converges to the same point in the topology \(\mathcal{C}_2\) .
\item
  Let \(\mathfrak{U}\) be an ultrafilter on \(\mathbf{N}\) which is finer than the Fréchet filter, and let \(\mathbf{X} = \mathbf{N} \cup \{\omega\}\) be the space associated with \(\mathfrak{U}\) (§ 6, no. 5). Show that a sequence \((x_n)\) which has an infinite number of distinct terms is never convergent in \(\mathbf{X}\) .
\item
  Let X, \(X'\) be two topological spaces and \(f: X \to X'\) a mapping which is continuous at a point \(x_{0} \in X\) . Show that if B is any filter base on X which has \(x_{0}\) as a cluster point, then \(f(x_{0})\) is a cluster point of the filter base \(f(\mathfrak{B})\) on \(X'\) .
\item
  \begin{enumerate}
  \def\labelenumii{\alph{enumii})}
  \tightlist
  \item
    Let X, Y be two topological spaces, F a filter on X having a cluster point a, and G a filter on Y having a cluster point b. Show that \((a, b)\) is a cluster point of the product filter \(F \times G\) .
  \end{enumerate}
\end{enumerate}

\begin{enumerate}
\def\labelenumi{\alph{enumi})}
\setcounter{enumi}{1}
\tightlist
\item
  In the product space \(Q^{2}\) , give an example of a sequence \((x_{n}, y_{n})\) which has no cluster points although each of the sequences \((x_{n})\) , \((y_{n})\) has a cluster point in Q.
\end{enumerate}

\begin{enumerate}
\def\labelenumi{\arabic{enumi})}
\setcounter{enumi}{4}
\item
  Let X, Y be two topological spaces, G the graph in \(X \times Y\) of a mapping \(f: X \to Y\) . Show that for each \(x \in X\) the set of cluster points of f at x is the section \(\overline{\mathbf{G}}(x)\) of \(\overline{\mathbf{G}}\) (the closure of G in \(X \times Y\) ) at x.
\item
  Let \((\mathbf{X}_t)_{t\in I}\) be an uncountable family of discrete spaces, each of which has at least two points. On the product set consider the topology
\end{enumerate}

\[
\mathbf {X} = \prod_ {\iota \in \mathbf {I}} \mathbf {X} _ {\iota},
\]

which is generated by all products \(\prod_{i\in I}M_i\) , where \(M_i\subset X_i\) for all \(i\in I\) and \(M_i = X_i\) for all but a countable number of indices (cf.~§ 4, Exercise 9). Show that every countable intersection of open sets is open in \(X\) , and deduce that no point of \(x\) can have a countable fundamental system of neighbourhoods.

\begin{enumerate}
\def\labelenumi{\arabic{enumi})}
\setcounter{enumi}{6}
\tightlist
\item
  A subset A of a topological space X is said to be primitive if it is the set of limit points of an ultrafilter on X. Let \(Y_{A}\) denote the set of ultrafilters on X whose set of limit points is A.
\end{enumerate}

\begin{enumerate}
\def\labelenumi{\alph{enumi})}
\item
  If A is a primitive subset of X, show that every open subset of X which meets A belongs to all the ultrafilters \(U \in Y_{A}\) .
\item
  Let A, B be two distinct primitive subsets of X. Show that there is an ultrafilter \(U \in Y_{A}\) , an ultrafilter \(B \in Y_{B}\) and a subset M of X such that \(M \in U\) and \(CM \in B\) .
\item
  Let \(f: \mathbf{X} \to \mathbf{Y}\) be a continuous mapping. Show that the image under \(f\) of any primitive subset of \(\mathbf{X}\) is contained in a primitive subset of \(\mathbf{Y}\) .
\item
  If \(x \in \mathbf{X}\) and \(\{x\}\) is closed in \(\mathbf{X}\) , show that \(\{x\}\) is a primitive subset of \(\mathbf{X}\) .
\end{enumerate}

\BourbakiExerciseGroup

\begin{enumerate}
\def\labelenumi{\arabic{enumi})}
\tightlist
\item
  \begin{enumerate}
  \def\labelenumii{\alph{enumii})}
  \tightlist
  \item
    Let X be a topological space. Show that the following three statements are equivalent:
  \end{enumerate}
\end{enumerate}

\begin{enumerate}
\def\labelenumi{(\Alph{enumi})}
\setcounter{enumi}{16}
\tightlist
\item
  If \(x, y\) are any two distinct points of \(X\) , there is a neighbourhood of \(x\) which does not contain \(y\) .
\end{enumerate}

(Q') Every subset of X which consists of a single point is closed in X. (Q'\,') For every \(x \in X\) the intersection of the neighbourhoods of x consists of the point x alone.

A space X is said to be accessible if it satisfies these conditions.

\begin{enumerate}
\def\labelenumi{\alph{enumi})}
\setcounter{enumi}{1}
\tightlist
\item
  An ordered set X with the right topology (with respect to which X is a Kolmogoroff space {[}§ 1, Exercise 2){]} is accessible if and only if no two distinct elements of X are comparable, and then the right topology on X is the discrete topology.
\end{enumerate}

\begin{enumerate}
\def\labelenumi{\arabic{enumi})}
\setcounter{enumi}{1}
\tightlist
\item
  \begin{enumerate}
  \def\labelenumii{\alph{enumii})}
  \tightlist
  \item
    If a topology on a set X has a finite subbase, and if X is accessible with respect to this topology, then X is finite and the topology is discrete.
  \end{enumerate}
\end{enumerate}

\begin{enumerate}
\def\labelenumi{\alph{enumi})}
\setcounter{enumi}{1}
\tightlist
\item
  If, in an accessible space X, every intersection of open sets is open, then X is discrete (cf.~§ 7, Exercise 6).
\end{enumerate}

\begin{enumerate}
\def\labelenumi{\arabic{enumi})}
\setcounter{enumi}{2}
\tightlist
\item
  \begin{enumerate}
  \def\labelenumii{\alph{enumii})}
  \tightlist
  \item
    Let X be an accessible space, A a subset of X, and x a point of \(\overline{A}\) which does not belong to A. Show that every neighbourhood of x contains an infinite number of points of A.
  \end{enumerate}
\end{enumerate}

\begin{enumerate}
\def\labelenumi{\alph{enumi})}
\setcounter{enumi}{1}
\item
  Let X be an accessible space and let A be a subset of X. Show that the set of points \(x \in \overline{A}\) such that every neighbourhood of x contains a point of A other than x is closed in X.
\item
  Let X be an accessible space. Show that the intersection of all neighbourhoods of any subset A of X is equal to A.
\end{enumerate}

\begin{enumerate}
\def\labelenumi{\arabic{enumi})}
\setcounter{enumi}{3}
\tightlist
\item
  Show that every subspace of a Kolmogoroff space (resp. accessible space) is a Kolmogoroff space (resp. accessible space); and that a topology which is finer than a topology of a Kolmogoroff space (resp. accessible space) is the topology of a Kolmogoroff space (resp. accessible space).
\end{enumerate}

¶ 5) a) Let X be an infinite set. Show that all the topologies \(\mathcal{G}_{\mathfrak{c}}\) on X (§ 2, Exercise 8) such that \(1 < c \leqslant\) card (X) are topologies of an accessible space but are not Hausdorff.

\begin{enumerate}
\def\labelenumi{\alph{enumi})}
\setcounter{enumi}{1}
\tightlist
\item
  Show that the intersection of all the Hausdorff topologies on X is the non-Hausdorff topology \(\mathfrak{G}_{\mathrm{card}(\mathbf{N})}\) , which is also the coarsest topology of an accessible space on X. (Use Exercise 8 of §2, and remark that there exist Hausdorff topologies on X in which there are countably infinite subsets of X which are not closed).
\end{enumerate}

\begin{enumerate}
\def\labelenumi{\arabic{enumi})}
\setcounter{enumi}{5}
\tightlist
\item
  \begin{enumerate}
  \def\labelenumii{\alph{enumii})}
  \tightlist
  \item
    Let X be a Hausdorff space and D a dense subset of X. Show that \(\text{card}(X) \leqslant 2^{2^{\text{card}(D)}}\) (observe that every point of X is the limit of a filter base on D). Show that this upper bound for card (X) cannot be improved (§ 4, Exercise 5 b)) and deduce that on an infinite set A the set of ultrafilters and the set of filters are each equipotent to \(\mathfrak{P}(\mathfrak{P}(A))\) .
  \end{enumerate}
\end{enumerate}

\begin{enumerate}
\def\labelenumi{\alph{enumi})}
\setcounter{enumi}{1}
\tightlist
\item
  Let K be the discrete space consisting of the two numbers o and 1. Let A be an infinite set. Deduce from a) that if card (A) ≥ 2 \(^{2\text{card(N)}}\) , then the product space K \(^{\text{A}}\) satisfies the condition (D \(_{\text{IV}}\) ) of Exercise 7 of § 1, but satisfies neither (D \(_{\text{II}}\) ) nor (D \(_{\text{III}}\) ). (§ 4, Exercise 4).
\end{enumerate}

* 7) Let X be the interval \([-1, 1]\) in R. Consider the following equivalence relation S on X: if \(x \neq \pm 1\) , the equivalence class of \(x\) consists of \(x\) and \(-x\) ; the equivalence class of \(1\) (resp. \(-1\) ) consists of \(1\) (resp. \(-1\) ) alone. Show that S is open and that the quotient space X/S is accessible but not Hausdorff. Show also that every point of X/S has a neighbourhood in X/S which is a regular subspace. *

\begin{enumerate}
\def\labelenumi{\arabic{enumi})}
\setcounter{enumi}{7}
\tightlist
\item
  \begin{enumerate}
  \def\labelenumii{\alph{enumii})}
  \tightlist
  \item
    Let X be a Hausdorff (resp. regular) space which is the sum of a family \((\mathbf{X}_{t})\) of subspaces, and let X/R be the quotient space obtained by pasting together the \(X_{t}\) along closed subsets \(A_{tx}\) by means of homeomorphisms \(h_{x_{t}}\) (§ 3, no. 4); suppose also that for each index t there is only a finite number of indices x such that \(A_{tx} \neq \emptyset\) . Under these conditions, show that X/R is Hausdorff (resp. regular).
  \end{enumerate}
\end{enumerate}

* b) Show that the space X/S of Exercise 7 can be obtained by pasting together two regular subspaces along two open sets. *

\begin{enumerate}
\def\labelenumi{\arabic{enumi})}
\setcounter{enumi}{8}
\tightlist
\item
  Let \(\Gamma\) be a finite group of homeomorphisms of a Hausdorff (resp. regular) space X, and let R be the equivalence relation ``there exists \(\sigma \in \Gamma\) such that \(y = \sigma(x)\)'' between \(x\) and \(y\) in X. Show that the quotient space X/R is Hausdorff (resp. regular).
\end{enumerate}

* 10) Let X be the subspace of R which is the complement of the set of points \(\mathbf{I} / n\) , where \(n\) is a rational integer other than \(\mathbf{o}, \mathbf{i}\) or \(-\mathbf{i}\) . Let S be the equivalence relation on X whose equivalence classes are (i) the set \(\{\mathbf{o}\}\) ; (ii) the set of all non-zero integers; (iii) all sets of the form \(\{x, \mathbf{i} / x\}\) where \(x \in \mathbf{X}\) and \(\mathbf{o} < |x| < \mathbf{i}\) .

\begin{enumerate}
\def\labelenumi{\alph{enumi})}
\item
  Show that the relation S on X is neither open nor closed.
\item
  Show that the graph of S is closed in \(X \times X\) , but that the quotient space X/S is not Hausdorff. *
\end{enumerate}

¶ * 11) Let X be a topological space. Let Z denote the product space \(R^{X}\) , and for each \(x \in X\) let \(Z_{x}\) be the subset of Z consisting of all \(u \in R^{X}\) such that \(u(x) \in Q\) and \(u(y) \notin Q\) for all \(y \in X\) other than x. Let Y be the subspace of the product \(X \times Z\) which is the union of the sets \(\{x\} \times Z_x\) as \(x\) runs through X. Show that Y is Hausdorff and that X is homeomorphic to a quotient space of Y. *

\begin{enumerate}
\def\labelenumi{\arabic{enumi})}
\setcounter{enumi}{11}
\tightlist
\item
  Let X be a topological space which has at least two points.
\end{enumerate}

\begin{enumerate}
\def\labelenumi{\alph{enumi})}
\item
  Show that the topologies \(\mathfrak{C}_{\Omega}\) and \(\mathfrak{C}_{\Phi}\) on \(\mathfrak{P}_0(\mathbf{X})\) (§ 2, Exercise 7) are not topologies of an accessible space.
\item
  Let \(\mathcal{G}_{\Theta}\) denote the least upper bound of the topologies \(\mathcal{G}_{\Omega}\) and \(\mathcal{G}_{\Phi}\) . Show that the set \(\mathfrak{F}(\mathbf{X})\) of non-empty closed subsets of \(\mathbf{X}\) , with the topology induced by \(\mathcal{G}_{\Theta}\) , is a Kolmogoroff space, and that if \(\mathbf{X}\) is accessible, so is \(\mathfrak{F}(\mathbf{X})\) .
\item
  Supposing that X is accessible, show that the space \(\mathfrak{F}(\mathbf{X})\) (with the topology induced by \(C_{\Theta}\) ) is Hausdorff if and only if X is regular.
\end{enumerate}

¶ 13) Let X be a set and let Γ be a group of permutations of X. a) Let F be a set of subsets of X. Let S be the set of all topologies on X in which all the sets of F are closed and every u ∈ Γ is a homeomorphism. Show that the set S has a least element C(Γ, F). If X is infinite, show that there exist sets F of subsets of X such that C(Γ, F) is accessible but not discrete (§ 2, Exercise 8).

\begin{enumerate}
\def\labelenumi{\alph{enumi})}
\setcounter{enumi}{1}
\tightlist
\item
  Let \(J(\Gamma)\) be the set of subsets \(A\) of \(X\) which have the following property: for each \(x \notin A\) there is a mapping \(f \in \Gamma\) such that \(f(x) \neq x\) which leaves fixed every element of \(A\) . Show that every Hausdorff topology on \(X\) , in which every mapping \(u \in \Gamma\) is a homeomorphism of \(X\) onto itself, is finer than \(\mathcal{C}_{J}(\Gamma) = \mathcal{C}(\Gamma, J(\Gamma))\) .
\end{enumerate}

Show that the topology \(\mathfrak{G}_{\mathrm{J}}(\Gamma)\) is Hausdorff if and only if the group \(\Gamma\) satisfies the following condition: given any two distinct elements \(x, y\) of \(\mathbf{X}\) there is a finite number of elements \(u_{k} \in \Gamma\) such that, if \(\mathbf{F}(u_{k})\) is the set of elements of \(\mathbf{X}\) which are invariant under \(u_{k}\) , the \(\mathbf{F}(u_{k})\) cover \(\mathbf{X}\) and none of the \(\mathbf{F}(u_{k})\) contains both \(x\) and \(y\) .

\begin{enumerate}
\def\labelenumi{\alph{enumi})}
\setcounter{enumi}{2}
\item
  Let \(\mathcal{C}_{0}\) denote the topology of the rational line \(\mathbf{Q}\) (*resp. the real line \(\mathbf{R}_{*}\) ) and let \(\Gamma\) be the group of all homeomorphisms of \(\mathbf{Q}\) (*resp. \(\mathbf{R}_{*}\) ) onto itself. Show that the topology \(\mathcal{C}_{\mathbf{J}}(\Gamma)\) is the same as \(\mathcal{C}_{0}\) .
\item
  Let X be the subspace of the rational line consisting of 0 and the points \(\iota/n\) , where n is an integer \(\geqslant\iota\) , and let \(\Gamma\) be the group of homeomorphisms of X. Show that the topology \(\mathcal{G}_{\mathrm{J}}(\Gamma)\) is not Hausdorff, and that the group of homeomorphisms of X with respect to this topology is \(\Gamma\) .
\item
  Let X be a topological space such that for each pair of distinct points x, y of X there exists a homeomorphism u of X onto itself such that \(u(x) = x\) and \(u(y) \neq y\) (*for example, the spaces \(R^{n}_{*}\) ). Let G be the group of homeomorphisms of X onto itself, and for each \(x \in X\) let \(S_{x}\) be the subgroup of G which leaves x fixed. Show that \(S_{x}\) is equal to its normalizer in G and that the intersection of all the subgroups
\end{enumerate}

\(S_{x}\) consists of the identity element alone; in particular, the centre of G consists of the identity element alone.

\begin{enumerate}
\def\labelenumi{\arabic{enumi})}
\setcounter{enumi}{13}
\tightlist
\item
  Show that the axiom \((\mathrm{O}_{\mathrm{III}})\) is equivalent to each of the following axioms:
\end{enumerate}

\((\mathrm{O}_{\mathrm{III}}^{\prime\prime})\) If F is any closed subset of X then the intersection of all the closed neighbourhoods of F is equal to F.

\((\mathrm{O}_{\mathrm{III}}^{\prime\prime})\) If F is any filter on X which converges to a point a, then the filter F on X which has the closures of the sets of F as a base also converges to a.

\begin{enumerate}
\def\labelenumi{\arabic{enumi})}
\setcounter{enumi}{14}
\tightlist
\item
  \begin{enumerate}
  \def\labelenumii{\alph{enumii})}
  \tightlist
  \item
    Show that every Kolmogoroff space which satisfies axiom (O \(_{III}\) ) is Hausdorff (and hence regular).
  \end{enumerate}
\end{enumerate}

\begin{enumerate}
\def\labelenumi{\alph{enumi})}
\setcounter{enumi}{1}
\tightlist
\item
  Construct a non-Hausdorff topology on a set of three elements which satisfies axiom \((\mathrm{O}_{\mathrm{III}})\) .
\end{enumerate}

\begin{enumerate}
\def\labelenumi{\arabic{enumi})}
\setcounter{enumi}{15}
\item
  Let \((\mathbf{X}_t)\) be an infinite family of topological spaces, and give the product set \(\mathbf{X} = \prod_{i} \mathbf{X}_i\) one of the topologies defined in Exercise 9 of §4. Show that if the \(\mathbf{X}_t\) are Kolmogoroff (resp. accessible, Hausdorff, regular) spaces then so is \(\mathbf{X}\) .
\item
  Show that the topologies \(\mathcal{C}_0(\mathbf{X})\) , \(\mathcal{C}_+(\mathbf{X})\) and \(\mathcal{C}_{-}(\mathbf{X})\) (§ 2, Exercise 5) on a linearly ordered set \(\mathbf{X}\) are regular.
\item
  Let \(X_{1}\) , \(X_{2}\) be two sets, let \(F_{i}\) be a filter on \(X_{i}\) (i = 1, 2), and let f be a mapping of \(X_{1} \times X_{2}\) into a regular space Y such that, for each
\end{enumerate}

\[
x _ {1} \in \mathrm{X} _ {1} \text { there   exists } \lim _ {x _ {2}, \mathfrak {F} _ {2}} f (x _ {1}, x _ {2}) = g (x _ {1}).
\]

Show that a point \(a \in Y\) is a cluster point of \(g\) with respect to the filter \(\mathfrak{F}_1\) if and only if, given any neighbourhood \(V\) of \(a\) in \(Y\) and any set \(\mathbf{A}_1 \in \mathfrak{F}_1\) , there exists \(x_1 \in \mathbf{A}_1\) and \(\mathbf{A}_2 \in \mathfrak{F}_2\) such that \(f(\{x_1\} \times \mathbf{A}_2) \subset V\) . ¶ 19) Let \(X\) be a Hausdorff space which is not regular (cf.~Exercise 20), and let \(a\) be a point of \(X\) such that there is a neighbourhood \(U\) of \(a\) which does not contain any closed neighbourhood of \(a\) . Let \(\mathfrak{F}\) be the neighbourhood filter of \(a\) , and let \(Y = \mathfrak{F} \cup \{\omega\}\) be the topological space associated with the section filter of \(\mathfrak{F}\) . Let \(Z\) be the set \(Y \times X\) with the following topology: for any point \((y, z) \neq (\omega, a)\) , the neighbourhoods of \((y, z)\) are the same as for the product topology, and the neighbourhoods of \((\omega, a)\) are the sets containing sets of the form

\[
(\{\omega \} \times \mathrm{V}) \cup ((\mathrm{Y} - \{\omega \}) \times \mathrm{X}),
\]

where \(\mathbf{V}\) is a neighbourhood of \(a\) in \(\mathbf{X}\) . Let \(\mathbf{A}\) be the subset of \(\mathbf{Z}\) which is the union of the sets \(\{\mathbf{V}\} \times \mathbf{V}\) where \(\mathbf{V} \in \mathfrak{F}\) . Let \(f\) be the mapping \((V, x) \rightarrow x\) of A into X. Show that f has a limit relative to A at each point of \(\overline{A}\) , but that there is no continuous mapping of \(\overline{A}\) into X which extends f.

¶ 20) a) An open subset U of a topological space X is said to be regular if \(\mathbf{U} = \alpha(u)\) (§ 1, Exercise 3); equivalently, U is regular if U is the interior of a closed set. X is said to be semi-regular (and its topology C is said to be semi-regular) if the regular open subsets of X form a base of C. Show that if C satisfies axiom \((\mathrm{O}_{\mathrm{III}})\) then C is semi-regular.

\begin{enumerate}
\def\labelenumi{\alph{enumi})}
\setcounter{enumi}{1}
\item
  Prove that the intersection of two regular open sets is a regular open set. Deduce that the regular open sets with respect to \(\mathcal{C}\) form a base of a topology \(\mathcal{C}^*\) on X, which is coarser than \(\mathcal{C}\) and semi-regular; \(\mathcal{C}^*\) is said to be the semi-regular topology associated with \(\mathcal{C}\) . Show that for each subset U of X which is open in \(\mathcal{C}\) , the closure of U with respect to \(\mathcal{C}^*\) is the same as its closure in \(\mathcal{C}\) , and that if X is accessible for \(\mathcal{C}\) the isolated points of X with respect to \(\mathcal{C}^*\) are the same as the isolated points of X with respect to \(\mathcal{C}\) . Show that \(\mathcal{C}^*\) is Hausdorff if and only if \(\mathcal{C}\) is {[}use Exercise 3 d) of §1{]}, and finally that if Y is a regular space then the continuous mappings of X into Y are the same in both topologies \(\mathcal{C}\) and \(\mathcal{C}^*\) .
\item
  Let X be a semi-regular space and let \(C_{0}\) be its topology. Show that a topology C on X is such that \(C^{*} = C_{0}\) if and only if there is a set M of dense subsets of X (in the topology \(C_{0}\) ) such that every finite intersection of sets of M belongs to M, and such that the topology C is generated by the union of M and the set of subsets of X which are open in the topology \(C_{0}\) . {[}Consider the dense open subsets (in the topology C) and notice that every open set in C is the intersection of a dense open set (in the topology C) and an open set (in the topology \(C_{0}\) ).{]} Hence construct examples of non-regular Hausdorff topologies which are finer than a regular topology (*and even finer than a topology of a compact metrizable space*).
\end{enumerate}

\begin{enumerate}
\def\labelenumi{\arabic{enumi})}
\setcounter{enumi}{20}
\tightlist
\item
  Let X be a Hausdorff space with no isolated points, and let \(C_{0}\) be the topology of X. Let M denote the set of complements of countable subsets A of X such that \(\overline{A}\) has only a finite number of non-isolated points. Show that if C is the topology on X generated by the union of M and the set of subsets which are open in \(C_{0}\) , then \(C^{*} = C_{0}\) (Exercise 20) and that every sequence which converges with respect to C has only a finite number of distinct terms (which implies that, with respect to C, no point of X has a countable fundamental system of neighbourhoods, although X itself may be countable and therefore every point of X is then the intersection of a countable family of neighbourhoods of this point).
\end{enumerate}

¶ 22) a) With the notation of Exercise 20 c), show that the set \(\mathrm{E}(\mathfrak{C}_{0})\) of topologies C such that \(C^{*}=C_{0}\) is inductive with respect to the relation ``C is coarser than \(C'\) ''. If \(C_{0}\) is a semi-regular topology on X, a maximal element of \(\mathrm{E}(\mathfrak{C}_{0})\) is called a submaximal topology, and the set X with such a topology is called a submaximal space. Thus every topology of \(\mathrm{E}(\mathfrak{C}_{0})\) is coarser than some submaximal topology.

\begin{enumerate}
\def\labelenumi{\alph{enumi})}
\setcounter{enumi}{1}
\item
  A topology \(\mathfrak{C}\) on \(\mathbf{X}\) is submaximal if and only if every subset of \(\mathbf{X}\) which is dense in the topology \(\mathfrak{C}\) is open in \(\mathfrak{C}\) . A non-empty submaximal space is therefore not solvable (§ 2, Exercise 2).
\item
  Show that every subspace of a submaximal space is locally closed and submaximal (consider first the open subspaces and the closed subspaces).
\item
  Show that the space associated with a filter which is finer than the filter of complements of finite subsets of an infinite set is a submaximal space.
\item
  If M is a subset of a submaximal space X, show that the subspace \(\overline{M} - \mathring{M} = \operatorname{Fr}(M)\) is discrete.
\item
  Conversely, let X be a topological space which has a dense open subset U such that U is submaximal and the topology of X is U-maximal (§ 3, Exercise 11). Then X is submaximal.
\end{enumerate}

\begin{enumerate}
\def\labelenumi{\arabic{enumi})}
\setcounter{enumi}{22}
\item
  Let X be a topological space and let A be a dense subset of X such that the topology of X is A-maximal (§ 3, Exercise 11). Let f be a continuous mapping of A into a topological space Y, such that at every point of X --- A the set of cluster points of f relative to A is not empty (*this will be the case, for example, if X is not empty and Y is quasi-compact*). Show that f can be extended by continuity to the whole of X.
\item
  In the interval {]}0,1{[} of the rational line, let A be the set of rational numbers of the form \(k/2^n\) , and B the set of rational numbers of the form \(k/3^n\) ( \(k\) an integer). Let X be the space obtained by giving {]}0,1{[} the topology generated by the open intervals {]}α,β{[} contained in X, and the sets A and B. Show that X is Hausdorff and semi-regular (Exercise 20) but not regular.
\end{enumerate}

¶ 25 a) On a set X, the set of regular topologies and the set of regular topologies with no isolated points are inductive sets with respect to the relation ``C is coarser than \(C'\)''.

\begin{enumerate}
\def\labelenumi{\alph{enumi})}
\setcounter{enumi}{1}
\item
  A topology \(\mathcal{G}\) on a set X is said to be ultraregular if it is maximal in the set of topologies on X which are regular and have no isolated points; hence given any regular topology on X with no isolated points, there is a finer ultraregular topology. A space is said to be ultraregular if its topology is ultraregular. Show that a regular topology \(\mathcal{G}\) with no isolated points is ultraregular if and only if, whenever A and B are complementary subsets of X which have no isolated points, then A and B are open sets. In particular, an ultraregular space is not solvable (§ 2, Exercise 11). (To prove that the condition is sufficient, observe that if \(\mathfrak{C}'\) is a regular topology with no isolated points and is finer than \(\mathfrak{G}\) , and if U is a non-empty open set with respect to \(\mathfrak{C}'\) and \(\overline{\mathrm{U}}\) is its closure in the topology \(\mathfrak{C}'\) , then \(\overline{\mathrm{U}}\) and \(\mathrm{X} - \overline{\mathrm{U}}\) have no isolated points with respect to the topology \(\mathfrak{C}'\) ).
\item
  In an ultraregular space, the closure of every set which has no isolated points is open, and the interior of a dense set is dense; the intersection of two dense sets is therefore dense. Deduce that if \(C_{0}\) is an ultraregular topology, then amongst all the topologies C such that \(C^{*} = C_{0}\) (Exercise 20) there is one (necessarily submaximal) which is finer than all the others.
\end{enumerate}

* 26) Let \((\theta_{n})\) be a sequence of irrational numbers which is dense in the interval \(\mathbf{I} = [\mathbf{o},\mathbf{i}]\) of \(\mathbf{R}\) . Let \(\mathbf{I}_n\) be the quotient space of \(\mathbf{I}\) obtained by identifying all the points \(0, \theta_1, \ldots, \theta_n\) , and let \(\varphi_n\) be the canonical mapping \(\mathbf{I} \to \mathbf{I}_n\) . Let \(\mathbf{X}\) be the set of rational numbers contained in \(\mathbf{I}\) ; then the restriction of \(\varphi_n\) to \(\mathbf{X}\) is a bijection of \(\mathbf{X}\) onto \(\varphi_n(\mathbf{X})\) . Let \(\mathcal{G}_n\) be the inverse image under this bijection of the topology induced on \(\varphi_n(\mathbf{X})\) by the topology of \(\mathbf{I}_n\) . Show that the greatest lower bound of the decreasing sequence \((\mathcal{G}_n)\) of Hausdorff topologies on \(\mathbf{X}\) is not a Hausdorff topology.*

¶ 27) Let X be a topological space. Show that there exists a Kolmogoroff (resp. accessible, Hausdorff, regular) space Y and a continuous mapping \(\varphi: X \to Y\) which has the following universal property: for every continuous mapping \(f: X \to Z\) of X into a Kolmogoroff (resp. accessible, Hausdorff, regular) space Z, there is a unique continuous mapping \(g: Y \to Z\) such that \(f = g \circ \varphi\) (Set Theory, Chapter IV, § 3, no. 1). Prove that:

\begin{enumerate}
\def\labelenumi{\alph{enumi})}
\item
  For Kolmogoroff spaces, we may take Y to be the quotient space of X by the equivalence relation \(\{\bar{x}\}=\{\bar{y}\}\) .
\item
  For accessible spaces, we may take Y to be the quotient space of X by the following equivalence relation R: the graph of R is the intersection of all graphs G of equivalences on X such that \(\mathrm{G}(x)\) is closed in x for each \(x \in X\) ; the graph of R then contains the graph of the relation \(\{\bar{x}\} \cap \{\bar{y}\} \neq \emptyset\) between x and y.
\item
  For Hausdorff spaces and regular spaces, show that conditions \((\mathrm{CU}_{\mathbf{I}})\) to \((\mathrm{CU}_{\mathbf{III}})\) of Set Theory, Chapter IV, § 3, no. 2, are satisfied, by using in particular Exercise 6 a). Give examples in which X is not regular and the mapping \(\varphi: X \to Y\) of X into the corresponding universal regular space is bijective {[}cf.~Exercise 20 b){]}.
\end{enumerate}

\begin{enumerate}
\def\labelenumi{\alph{enumi})}
\setcounter{enumi}{3}
\tightlist
\item
  If X is a topological space which satisfies axiom \((\mathrm{O}_{\mathrm{III}})\) , show that the universal Kolmogoroff space corresponding to X is canonically homeomorphic to the universal regular space corresponding to X.
\end{enumerate}

\begin{enumerate}
\def\labelenumi{\arabic{enumi})}
\setcounter{enumi}{27}
\tightlist
\item
  Let X be a Hausdorff space which is not regular, let F be a closed subset of X and let \(a \notin F\) be a point of X such that every neighbourhood of a meets every neighbourhood of F. Let R be the equivalence relation on X obtained by identifying all the points of F. Show that R is closed and that the graph of R is closed in \(X \times X\) , but that R is not Hausdorff.
\end{enumerate}

¶ 29) a) Let R be a Hausdorff equivalence relation on a non-empty topological space X, and let \(\mathfrak{F}(X)\) be the space of non-empty closed subsets of X, with the topology induced by \(\mathfrak{G}_{\Theta}\) (Exercise 12). Show that every \(A \in \mathfrak{F}(X)\) which lies in the closure of X/R {[}considered as a subset of \(\mathfrak{F}(X)\) {]} is contained in some equivalence class of R (use reductio ad absurdum).

\begin{enumerate}
\def\labelenumi{\alph{enumi})}
\setcounter{enumi}{1}
\tightlist
\item
  Suppose in addition that X is regular and R is open. Show that X/R is then a closed subset of \(\mathfrak{F}(X)\) with respect to the topology \(\mathfrak{G}_{\Theta}\) (use Proposition 11 of § 7, no. 6).
\end{enumerate}

\BourbakiExerciseGroup

\begin{enumerate}
\def\labelenumi{\arabic{enumi})}
\item
  Let X be a topological space, S a subbase of the topology of X (§ 2, no. 3). Show that X is quasi-compact if every open covering of X formed by sets belonging to S contains a finite covering of X. (Observe that if an ultrafilter U on X were not convergent, then every \(x \in X\) would belong to a set of S not belonging to U; and use the corollary to Proposition 5 of § 6).
\item
  Show that a well-ordered set X with the topology \(\mathcal{C}_{-}(\mathbf{X})\) (§ 2, Exercise 5) is locally compact, and is compact if and only if X has a greatest element. Deduce that every set can be topologized in such a way that it becomes a compact space.
\item
  \begin{enumerate}
  \def\labelenumii{\alph{enumii})}
  \tightlist
  \item
    Let X be a Hausdorff space, and let A, B be two disjoint compact subsets of X. Show that there is a neighbourhood of A and a neighbourhood of B which have no common point.
  \end{enumerate}
\end{enumerate}

\begin{enumerate}
\def\labelenumi{\alph{enumi})}
\setcounter{enumi}{1}
\tightlist
\item
  Let X be a regular space and A a compact subset of X, and let B be a closed subset of X which does not meet A. Show that there is a neighbourhood of A and a neighbourhood of B which have no common point.
\end{enumerate}

\begin{enumerate}
\def\labelenumi{\arabic{enumi})}
\setcounter{enumi}{3}
\item
  Show that in the space X defined in Exercise 6 of § 7, which is regular (§ 8, Exercise 16), every compact subset is finite.
\item
  *a) Show that the non-Hausdorff accessible space Y = X/S defined in Exercise 7 of § 8 is such that every point of Y has a compact neighbourhood; that the images α, β of 1 and -1 in Y have compact neighbourhoods which are not closed; that the intersection of a compact neighbourhood of α and a compact neighbourhood of β is not quasi-compact, and that the union of these two neighbourhoods is not compact.*
\end{enumerate}

\begin{enumerate}
\def\labelenumi{\alph{enumi})}
\setcounter{enumi}{1}
\item
  Let X be the sum of N and an infinite set A. For each \(x \in X\) let \(\mathfrak{S}(x)\) be the filter base defined as follows: if \(x \in N\) , then \(\mathfrak{S}(x)\) consists of the single set \(\{x\}\) ; if \(x \in A\) and if \(S_{x,n}\) denotes the set whose elements are x and the integers \(\geqslant n\) , then \(\mathfrak{S}(x)\) consists of the sets \(S_{x,n}\) where n runs through N. Show that there is a topology on X for which \(\mathfrak{S}(x)\) is a fundamental system of neighbourhoods of x for every \(x \in X\) , and that X is accessible but not Hausdorff in this topology. Show also that X is not quasi-compact, but that X has dense compact subsets.
\item
  Show that two topologies on a set X, with respect to each of which X is accessible and quasi-compact, can be distinct but comparable {[}cf.~§ 8, Exercises 7 and 5 b){]}.
\end{enumerate}

\begin{enumerate}
\def\labelenumi{\arabic{enumi})}
\setcounter{enumi}{5}
\item
  Let X, Y be two topological spaces, and let A (resp. B) be a quasi-compact subset of X (resp. Y). Show that if U is any neighbourhood of \(A \times B\) in \(X \times Y\) then there is a neighbourhood V of A in X and a neighbourhood W of B in Y such that \(V \times W \subset U\) .
\item
  \begin{enumerate}
  \def\labelenumii{\alph{enumii})}
  \tightlist
  \item
    Give an example of an inverse system \((\mathbf{X}_{\alpha}, f_{\alpha\beta})\) of quasi-compact spaces whose inverse limit is empty (observe that in the space of Example 2 of no. 1, every subspace is quasi-compact).
  \end{enumerate}
\end{enumerate}

\begin{enumerate}
\def\labelenumi{\alph{enumi})}
\setcounter{enumi}{1}
\tightlist
\item
  Let \(p\) be a prime number. Let \(Z_{n}\) denote the set \(\mathbf{Z}\) of rational integers, topologized with the inverse image of the discrete topology under the canonical mapping of \(\mathbf{Z}\) onto \(\mathbf{Z} / p^{n}\mathbf{Z}\) . If \(m \leqslant n\) , let \(f_{mn}\) denote the identity mapping \(Z_{n} \to Z_{m}\) , which is continuous. Show that the spaces \(Z_{n}\) are quasi-compact, but that the inverse limit of the inverse system \((Z_{n}, f_{mn})\) is not quasi-compact (*).
\end{enumerate}

\begin{enumerate}
\def\labelenumi{\arabic{enumi})}
\setcounter{enumi}{7}
\tightlist
\item
  Let X be a topological space and let \((\mathrm{G}_{\alpha})_{\alpha\in\mathbb{A}}\) be a family of graphs of equivalence relations on X, whose index set A is directed; let \(R_{\alpha}\) be the relation \((x,y)\in\mathbf{G}_{\alpha}\) and let \(\varphi_{\alpha}\) be the canonical mapping \(X\to X/R_{\alpha}\) . Suppose that the relation \(R_{\beta}\) implies \(R_{\alpha}\) if \(\alpha\leqslant\beta\) and let \(\varphi_{\alpha\beta}\) be the canonical mapping \(X/R_{\beta}\to X/R_{\alpha}\) ; then \((X/R_{\alpha},\varphi_{\alpha\beta})\) is an inverse system of topological spaces. Let \(Y=\varprojlim X/R_{\alpha}\) .
\end{enumerate}

\begin{enumerate}
\def\labelenumi{\alph{enumi})}
\item
  Let \(g: \mathbf{X} \to \mathbf{Y}\) be the continuous mapping \(\varprojlim \varphi_{\alpha}\) . Show that \(g(x)\) is dense in \(\mathbf{Y}\) and that, for each \(x \in \mathbf{X}\) , \(\overrightarrow{g}(g(x))\) is the class of \(x\) with respect to the equivalence relation whose graph is \(\bigcap_{n} G_{\alpha}\) .
\item
  Suppose that all the \(R_{\alpha}\) are closed equivalence relations, and that for each \(x \in X\) and each neighbourhood V of x there is an index \(\alpha\) such that the class of x with respect to \(R_{\alpha}\) is contained in V. Show that under these conditions g is an open mapping of X onto \(g(\mathbf{X})\) .
\item
  Show that if the classes mod \(R_{\alpha}\) are closed and quasi-compact for each \(\alpha \in A\) , then g is surjective.
\end{enumerate}

\begin{enumerate}
\def\labelenumi{\arabic{enumi})}
\setcounter{enumi}{8}
\item
  Let X be an accessible space and \(\mathcal{R}\) a set of quasi-compact closed subsets of X, which contains all finite subsets of X and is such that the union of any two subsets of \(\mathcal{R}\) belongs to \(\mathcal{R}\) . Let \(X'\) be a set which is the sum of X and a set \(\{\omega\}\) consisting of a single point. Show that there exists a quasi-compact accessible topology on \(X'\) which induces the given topology on X, and which is such that the complements in \(X'\) of the sets of \(\mathcal{R}\) form a fundamental system of open neighbourhoods of \(\omega\) . Give examples of distinct topologies on \(X'\) corresponding to distinct sets \(\mathcal{R}\) .
\item
  Let X be a paracompact space. Show that if every open subspace of X is paracompact, then every subspace of X is paracompact. Deduce that every subspace of a locally compact space whose topology has a countable base, is paracompact {[}cf.~Chapter IX, § 4, Exercise 20 a){]}.
\end{enumerate}

¶ 11) Let \(\mathbf{X}_0\) be an uncountable well-ordered set which has a greatest element; let \(a\) be the smallest element of \(\mathbf{X}_0\) and \(b\) the smallest of the elements \(x \in \mathbf{X}_0\) such that the interval \([a, x[\) is uncountable. Let \(\mathbf{X}\) denote the interval \([a, b[\) with the topology \(\mathfrak{C}_{-}(\mathbf{X})\) .

\begin{enumerate}
\def\labelenumi{\alph{enumi})}
\item
  X is locally compact but not compact (Exercise 2) and every sequence in X has a cluster point (observe that every countable subset of X is bounded above).
\item
  Let \(f\) be a mapping of \(\mathbf{X}\) into itself such that \(f(x) < x\) for all sufficiently large \(x\) . Show that there exists an element \(c \in \mathbf{X}\) with the property that for each \(x \in \mathbf{X}\) there exists \(y \geqslant x\) such that \(f(y) \leqslant c\) . {[}Use reductio ad absurdum, by constructing a sequence \((z_n)\) of points of \(\mathbf{X}\) such that \(z_{n+1}\) is the smallest of the elements \(z' \in \mathbf{X}\) such that \(f(x) \geqslant z_n\) whenever \(x \geqslant z'\) {]}.
\item
  Deduce from b) that X is not paracompact.
\end{enumerate}

\begin{enumerate}
\def\labelenumi{\arabic{enumi})}
\setcounter{enumi}{11}
\tightlist
\item
  Let X be a topological space, \(\mathfrak{F}(X)\) the set of non-empty closed subsets of X, \(\mathfrak{K}(X)\) the set of non-empty quasi-compact subsets of X;
\end{enumerate}

endow each of these sets with the topology induced by the topology \(\mathfrak{C}_{\Omega}\) on \(\mathfrak{P}_{0}(\mathbf{X})\) (§ 2, Exercise 7).

\begin{enumerate}
\def\labelenumi{\alph{enumi})}
\tightlist
\item
  Show that if \(\mathfrak{B} \subset \mathfrak{F}(\mathbf{X})\) is quasi-compact with respect to \(\mathcal{C}_{\Omega}\) and if \(\mathbf{X}\) is regular, then the union of the sets \(\mathbf{M} \in \mathfrak{B}\) is closed in \(\mathbf{X}\) .\\
\item
  Show that if \(\mathfrak{B} \subset \mathfrak{K}(\mathbf{X})\) is quasi-compact with respect to \(\mathcal{C}_{\Omega}\) then the union of the sets \(\mathbf{M} \in \mathfrak{B}\) is quasi-compact.
\end{enumerate}

¶ 13) With the notation of Exercise 13, endow \(\mathfrak{F}(\mathbf{X})\) and \(\mathfrak{R}(\mathbf{X})\) with the topology induced by the topology \(\mathfrak{C}_{\Theta}\) on \(\mathfrak{P}_{0}(\mathbf{X})\) (§ 8, Exercise 12). a) Show that if X is quasi-compact, then so is \(\mathfrak{F}(\mathbf{X})\) . (For each open set U in X, let \(c(\mathrm{U})\) {[}resp. \(m(\mathrm{U})\) {]} be the set of all closed subsets M of X such that \(\mathrm{M} \subset \mathrm{U}\) (resp. \(\mathrm{M} \cap \mathrm{U} \neq \emptyset\) ); the \(c(\mathrm{U})\) and \(m(\mathrm{U})\) form a subbase of \(\mathfrak{C}_{\Theta}\) . Then use Exercise 1).

\begin{enumerate}
\def\labelenumi{\alph{enumi})}
\setcounter{enumi}{1}
\item
  Suppose that X is accessible, and let \(X'\) be the image of X under the mapping \(x \to \{x\}\) of X into \(\mathfrak{F}(X)\) . Let H be a subspace of \(\mathfrak{F}(X)\) which contains \(X'\) . Show that if H is quasi-compact then so is X. {[}If \((\mathrm{U}_{\alpha})\) is an open covering of X, remark that if \(B_{\alpha}\) denotes the set of all \(M \in \mathfrak{F}(X)\) such that \(M \cap U_{\alpha} \neq \emptyset\) , then \(\{\mathfrak{B}_{\alpha}\}\) is an open covering of \(\mathfrak{F}(X)]\) .
\item
  Show that if X is locally compact, then \(\mathfrak{K}(\mathbf{X})\) is open in \(\mathfrak{F}(\mathbf{X})\) (use Proposition 10 of § 7).
\item
  Suppose that X is accessible. Show that \(\mathcal{R}(X)\) is Hausdorff (resp. regular, compact, locally compact) if and only if X is. {[}For the first two assertions, use Exercise 12 of § 8 and Exercise 3 of § 9; for the last two assertions, use a) and b) and Exercise 12{]}.
\end{enumerate}

¶ 14) A topological space X is said to be a Lindelöf space if every open covering of X contains a countable covering of X. Every space with a countable base is a Lindelöf space, and so is every quasi-compact space.

\begin{enumerate}
\def\labelenumi{\alph{enumi})}
\item
  Every closed subspace of a Lindelöf space is a Lindelöf space. For every subspace of a Lindelöf space to be a Lindelöf space it is necessary and sufficient that every open subspace is Lindelöf (cf.~Exercise 15).
\item
  If \(f\) is any continuous mapping of a Lindelöf space \(\mathbf{X}\) into a topological space \(\mathbf{X}'\) , then the subspace \(f(\mathbf{X})\) of \(\mathbf{X}'\) is Lindelöf.
\item
  Every countable union of Lindelöf subspaces of a topological space is Lindelöf. In particular every countable space is Lindelöf.
\item
  A Lindelöf space X is quasi-compact if every sequence in X has a cluster point {[}show that axiom (C'') is then verified{]}.
\end{enumerate}

¶ 15) a) Let X be a topological space in which every open subspace is Lindelöf. Show that X satisfies condition (D \(_{III}\) ) of § 1, Exercise 7, and that if X is regular then every closed subset of X is a countable intersection of open sets {[}cf.~Exercise 23 c){]}.

\begin{enumerate}
\def\labelenumi{\alph{enumi})}
\setcounter{enumi}{1}
\item
  Let X be a topological space in which every open subspace is paracompact. Show that every open subspace of X is Lindelöf if and only if X satisfies condition \((D_{\mathrm{III}})\) of §1, Exercise 7.
\item
  Let X be a compact space. Show that every open subspace of X is Lindelöf if and only if (i) X satisfies condition \((\mathbf{D}_{\mathbf{III}})\) of §1, Exercise 7, and (ii) every closed subset of X is a countable intersection of open sets (use b) and Theorem 5 of no. 10).
\end{enumerate}

¶ 16) Let X be a topological space and A a subset of X. A point \(a \in X\) is said to be a condensation point of A if every neighbourhood of \(a\) contains an uncountable infinity of points of A.

\begin{enumerate}
\def\labelenumi{\alph{enumi})}
\item
  The set of condensation points of a subset A of X is closed. Give an example in which X has only one condensation point (cf.~no. 8, Theorem 4).
\item
  Suppose X is an accessible space in which every open subspace is Lindelöf (Exercise 15). Show that for every subset A of X, the set B of condensation points of A is perfect, and that A ∩ GB is countable.
\end{enumerate}

¶ 17) In a topological space X, a point \(a\) is said to be totally adherent to a subset A of X if for each neighbourhood U of \(a\) we have card (A) = card (A ∩ U). Show that X is quasi-compact if and only if for each infinite subset A of X there is a point of X which is totally adherent to A. (If X is not quasi-compact, show that there is an initial ordinal \(\omega_{\alpha}\) (Set Theory, Chapter III, § 6, Exercise 10) and an open covering (U \(_{\xi}\) ) of X, where \(\xi\) runs through the set of ordinals \(<\omega_{\alpha}\) , such that for each index \(\xi<\omega_{\alpha}\) the complement A \(_{\xi}\) in X of the union of the U \(_{\eta}\) with indices \(\eta<\xi\) has a cardinal equal to N \(_{\alpha}\) ; deduce that it is possible to define a family ( \(x_{\xi}\) ) ( \(\xi<\omega_{\alpha}\) ) of points of X by transfinite induction such that \(x_{\xi}\in\mathrm{A}_{\xi}\) for each \(\xi\) and \(x_{\xi}\neq x_{\eta}\) whenever \(<\eta\xi\) ).

¶ 18) A Hausdorff space X is said to be absolutely closed if it has the following property: for every homeomorphism f of X onto a subspace of a Hausdorff space X', f(X) is closed in X'. Every compact space is absolutely closed.

\begin{enumerate}
\def\labelenumi{\alph{enumi})}
\tightlist
\item
  Show that the following conditions are equivalent:
\end{enumerate}

(AF) The space X is absolutely closed.

(AF') Every filter base on X which consists of open sets has at least one cluster point.

(AF'') Every family of closed subsets of X whose intersection is empty contains a finite subfamily such that the interiors of the subsets in this subfamily have an empty intersection.

(AF''') Every open covering of X contains a finite subfamily whose closures cover X.

\begin{enumerate}
\def\labelenumi{\alph{enumi})}
\setcounter{enumi}{1}
\item
  Let X be a Hausdorff space and U a dense open subset of X. Suppose that every filter base on U which consists of open sets has at least one cluster point in X. Show that under these conditions X is absolutely closed. In particular, if A is an open subset of an absolutely closed space X, then \(\overline{A}\) is absolutely closed.
\item
  Let X be an absolutely closed space and let f be a continuous mapping of X into a Hausdorff space \(X'\) . Show that \(f(\mathbf{X})\) is an absolutely closed subspace of \(X'\) .
\item
  Show that the product of two absolutely closed spaces is absolutely closed (argue as in Proposition 17 of no. 10).
\end{enumerate}

¶ 19) A Hausdorff space X is said to be minimal if its topology C is such that every topology on X which is Hausdorff and coarser than C necessarily coincides with C; then C must be semi-regular (§ 8, Exercise 20). Show that a Hausdorff space X is minimal if and only if every filter base on X, which consists of open sets and has at most one cluster point, is convergent. It comes to the same thing to say that X is absolutely closed (Exercise 18) and that every filter base on X, which consists of open sets and has a single cluster point, converges to this point.

¶ 20) A topological space X is said to be completely Hausdorff (and the topology of X is said to be completely Hausdorff) if each pair of distinct points of X has disjoint closed neighbourhoods.

\begin{enumerate}
\def\labelenumi{\alph{enumi})}
\item
  Every regular space is completely Hausdorff. Every subspace of a completely Hausdorff space is completely Hausdorff. Every product of completely Hausdorff spaces is completely Hausdorff. Every topology which is finer than a completely Hausdorff topology is completely Hausdorff.
\item
  Every minimal completely Hausdorff space X (Exercise 19) is compact. (By reductio ad absurdum: consider an ultrafilter U on X and the filter base formed by the open sets which belong to U).
\end{enumerate}

\(^{*}c)\) Let \(\theta\) be an irrational number \textgreater0. For every point \((x,y)\in\mathbf{Q}^{2}\) and every \(\alpha>0\) let \(\mathrm{B}_{\alpha}(x,y)\) be the set consisting of \((x,y)\) and the points \((z,0)\in\mathbf{Q}^{2}\) such that \(|z-(x+\theta y)|<\alpha\) or \(|z-(x-\theta y)|<\alpha\) . Let \(\mathfrak{B}(x,y)\) be the set of \(\mathrm{B}_{\alpha}(x,y)\) where \(\alpha\) runs through the set of real numbers \textgreater0. Show that there is a Hausdorff topology \(\mathcal{C}\) on \(\mathbf{Q}^{2}\) such that, for each \((x,y)\) , \(\mathfrak{B}(x,y)\) is a fundamental system of neighbourhoods of \((x,y)\) for \(\mathcal{C}\) . Show that, given any two points a, b of \(\mathbf{Q}^{2}\) , every closed neighbourhood of a meets every closed neighbourhood of b in the topology \(\mathcal{C}.*\)

¶ 21) a) Let X be a Hausdorff space and let \(\mathfrak{G}\) be the topology of X. Show that X is absolutely closed (Exercise 18) if and only if the semi-regular topology \(\mathfrak{G}^*\) associated with \(\mathfrak{G}\) ( \(\S\) 8, Exercise 20) is the topology of a minimal space (Exercise 19).

\begin{enumerate}
\def\labelenumi{\alph{enumi})}
\setcounter{enumi}{1}
\tightlist
\item
  If \(\mathfrak{C}\) is completely Hausdorff (Exercise 20), then so is \(\mathfrak{C}_{*}\) (use Exercise 20 of § 8). Deduce that an absolutely closed regular space is compact.
\end{enumerate}

¶ 22) a) Let X be a Hausdorff space and let A be a dense subset of X. Show that if every filter on A has at least one cluster point in X, then X is absolutely closed {[}verify axiom (AF') of Exercise 18{]}.

\begin{enumerate}
\def\labelenumi{\alph{enumi})}
\setcounter{enumi}{1}
\tightlist
\item
  Let \(X_{0}\) be a compact space which has a non-open dense subset A, and let \(C_{0}\) be the topology of \(X_{0}\) . Let C be the topology on \(X_{0}\) generated by A and the subsets of \(X_{0}\) which are open in \(C_{0}\) , so that \(C^{*} = C_{0}\) (§ 8, Exercise 20). If X denotes the non-compact absolutely closed space obtained by endowing the set \(X_{0}\) with the topology C (Exercise 21) show that every filter on A has a cluster point in X.
\end{enumerate}

¶ 23) a) Let X be a countable absolutely closed space. Show that the set of isolated points of X is dense in X. {[}Use reductio ad absurdum, by arranging the points of X in a sequence, and constructing by induction a countable open covering (U \(_{n}\) ) of X such that no finite union of the \(\overline{U}_{n}\) covers X{]}.

*b) Give an example of a completely Hausdorff absolutely closed topological space which is not a Lindelöf space {[}in the product \([0,1]\times[0,1]\) ; consider the complements of the sets \(A\times\{o\}\) , where A runs through the set of all subsets of \([0,1]\) , and apply the general method of § 8, Exercise 20 c){]}. This space satisfies condition \((D_{II})\) of § 1, Exercise 1, but does not satisfy condition \((D_{III})\) .

\begin{enumerate}
\def\labelenumi{\alph{enumi})}
\setcounter{enumi}{2}
\tightlist
\item
  Let \(X_{0}\) be a compact space with no isolated points, let M be the set of complements of countable subsets of \(X_{0}\) {[}so that the sets of M are dense in \(X_{0}\) , cf.~a){]}. Let C be the topology generated by the open sets of \(X_{0}\) and the sets of M. Show that the space X obtained by giving \(X_{0}\) the topology C is absolutely closed and completely Hausdorff and is a Lindelöf space which does not satisfy condition ( \(D_{II}\) ) of §1, Exercise 7; show also that no point of X has a countable fundamental system of neighbourhoods and that no sequence in X, all of whose terms are distinct, has a cluster point. If every open subspace of \(X_{0}\) is Lindelöf {[}cf.~Exercise 16 c){]}, show that the same is true of X and hence in particular that X satisfies ( \(D_{III}\) ) {[}Exercise 16 a){]}. This is the case when \(X_{0} = [0,1]\) . Show however that in this case there are countable closed subsets of X which are not countable intersections of open sets of X (cf.~Chapter IX, §5, no. 3). *
\end{enumerate}

¶ 24) a) Let X be a Hausdorff space, and let Γ be the set of filters on X which have a base consisting of open sets and which have no cluster point in X. Show that the set Γ is inductive if it is not empty (i.e., if X is not absolutely closed).

\begin{enumerate}
\def\labelenumi{\alph{enumi})}
\setcounter{enumi}{1}
\tightlist
\item
  Let \(\Omega\) be the set of maximal filters in \(\Gamma\) , and let \(X'\) be the set which is the sum of \(X\) and \(\Omega\) . For each \(x \in X\) let \(\mathfrak{B}(x)\) denote the set of neighbourhoods of \(x\) in \(X\) , and for each \(\omega \in \Omega\) let \(\mathfrak{B}(\omega)\) denote the set of subsets of \(X'\) of the form \(\{\omega\} \cup M\) , where \(M\) belongs to the filter \(\omega\) . Then there is a topology on \(X'\) in which \(\mathfrak{B}(x')\) is a fundamental system of neighbourhoods of \(x'\) for each \(x' \in X'\) . Show that, with this topology, \(X'\) is absolutely closed and has the following universal property: given any open continuous mapping \(f: X \to Y\) of \(X\) into an absolutely closed space \(Y\) , there is a unique continuous mapping \(g: X' \to Y\) which extends \(f\) (consider the images under \(f\) of the filters which belong to \(\Omega\) , and observe that in an absolutely closed space, a filter which is maximal in the set of filters which have a base consisting of open sets must be convergent). In particular, the topology of \(X'\) is X-maximal ( \(\S 3\) , Exercise 11). Deduce that there are absolutely closed spaces whose topology is quasi-maximal ( \(\S 2\) , Exercise 6).
\end{enumerate}

¶ 25) a) Let X be an accessible space, and let \(\Phi\) be the set of filters on X which have a base consisting of closed sets. Show that the set \(\Phi\) is inductive, and let \(X''\) be the set of maximal filters in \(\Phi\) .

\begin{enumerate}
\def\labelenumi{\alph{enumi})}
\setcounter{enumi}{1}
\tightlist
\item
  For each open subset U of X, let \(U^{*}\) be the set of filters \(F \in X''\) such that \(U \in F\) . If U, V are two open subsets of X, then
\end{enumerate}

\[
(\mathrm{U} \cap \mathrm{V}) ^ {*} = \mathrm{U} ^ {*} \cap \mathrm{V} ^ {*} \quad \text { and } \quad (\mathrm{U} \cup \mathrm{V}) ^ {*} = \mathrm{U} ^ {*} \cup \mathrm{V} ^ {*}
\]

(remark that if \(\mathfrak{F} \in \mathbf{X}''\) is such that \(\mathbf{U} \notin \mathfrak{F}\) , then \(\mathbf{X} = \mathbf{U} \in \mathfrak{F}\) ). Show that the sets \(\mathbf{U}^*\) form a base of a topology on \(\mathbf{X}''\) and that in this topology \(\mathbf{X}''\) is quasi-compact {[}verify axiom \((\mathbf{C}''')\) , using the remark above{]} and accessible {[}if \(\mathfrak{F}, \mathfrak{G}\) are two distinct elements of \(\mathbf{X}''\) , show that there exist two disjoint closed subsets \(\mathbf{A}\) and \(\mathbf{B}\) in \(\mathbf{X}\) such that \(\mathbf{A} \in \mathfrak{F}\) and \(\mathbf{B} \in \mathfrak{G}\) {]}.

\begin{enumerate}
\def\labelenumi{\alph{enumi})}
\setcounter{enumi}{2}
\item
  For each \(x \in \mathbf{X}\) let \(\varphi(x)\) be the ultrafilter on \(\mathbf{X}\) which has the single set \(\{x\}\) as a base. Show that \(\varphi\) is a homeomorphism of \(\mathbf{X}\) onto a dense subset of \(\mathbf{X}''\) and that if \(\mathbf{X}\) is quasi-compact (and accessible), then \(\varphi(\mathbf{X}) = \mathbf{X}''\) .
\item
  Show that for each continuous mapping \(f\) of \(\mathbf{X}\) into a compact space \(\mathbf{Y}\) there is a unique continuous mapping \(g: \mathbf{X}'' \to \mathbf{Y}\) such that \(f = g \circ \varphi\) .
\item
  Let \(\tilde{\mathbf{X}}\) be the universal Hausdorff space associated with \(\mathbf{X}''\) (§ 8, Exercise 27); show that the canonical mapping \(\psi: \mathbf{X}'' \to \tilde{\mathbf{X}}\) is surjective and that \(\tilde{\mathbf{X}}\) is compact. Thus every continuous mapping of \(\mathbf{X}\) into a compact space \(\mathbf{Y}\) can be written uniquely in the form \(f = h \circ \psi \circ \varphi\) , where \(h: \tilde{\mathbf{X}} \to \mathbf{Y}\) is continuous.
\end{enumerate}

\begin{enumerate}
\def\labelenumi{\arabic{enumi})}
\setcounter{enumi}{25}
\tightlist
\item
  If X is a discrete space, show that the spaces \(X''\) and \(\tilde{X}\) defined in Exercise 25 are identical, and that if C is the topology of the space \(X'\) defined in Exercise 24, then \(X''\) can be identified with the space \(X'\) carrying the semi-regular topology \(C^*\) associated with C (§ 8, Exercise 20). Furthermore, the set \(X''\) can be canonically identified with the set of ultrafilters on X. The compact space \(X''\) is called the ultrafilter space of X.
\end{enumerate}

¶ 27) Show that in a compact space X every isodyne {[}§ 2, Exercise 11 d){]} infinite open subspace is solvable. {[}Using Proposition 1 of no. 2, show that there is a base B of the topology of U such that card (B) = card (U); then use Exercise 11 c) of § 2.{]} Deduce that every locally compact space with no isolated points is solvable {[}use Exercise 11 d) of § 2{]}; consequently such a space is not submaximal.

\begin{enumerate}
\def\labelenumi{\arabic{enumi})}
\setcounter{enumi}{27}
\tightlist
\item
  \begin{enumerate}
  \def\labelenumii{\alph{enumii})}
  \tightlist
  \item
    Let K be the discrete space \(\{0,1\}\) . If X is any compact space, let F be the set of mappings f of K into the set \(\mathfrak{F}(X)\) of closed subsets of X such that \(\mathbf{X}=f(0)\cup f(1)\) . Consider the compact space \(Y=K^{F}\) ; show that for each \(y=(y_{f})_{f\in\mathbb{F}}\) in Y, the intersection of the closed sets \(f(y_{f})\) of X consists of at most one point. The set Z of \(y\in Y\) for which this intersection is not empty is a closed subset of Y. If \(\varphi(y)\) is the unique point common to the \(f(y_{f})\) for \(y\in Z\) , show that \(\varphi\) is a continuous surjection of Z onto X.
  \end{enumerate}
\end{enumerate}

\begin{enumerate}
\def\labelenumi{\alph{enumi})}
\setcounter{enumi}{1}
\tightlist
\item
  Give an example of a compact space X such that there is no surjective continuous mapping of a compact space of the form \(K^{A}\) onto X {[}use Exercise 4 b) of § 4 and Exercise 26 of § 9{]}.
\end{enumerate}

§ 29) A topological space X is said to be locally quasi-compact if each point of X has a fundamental system of quasi-compact neighbourhoods. a) Show that in a locally quasi-compact space, every quasi-compact subset has a fundamental system of quasi-compact neighbourhoods.

\begin{enumerate}
\def\labelenumi{\alph{enumi})}
\setcounter{enumi}{1}
\tightlist
\item
  Show that every locally closed subspace of a locally quasi-compact space is locally quasi-compact.
\end{enumerate}

\(^{*}c)\) On the closed Euclidean disc B: \(||x|| \leqslant 1\) in the real plane \(R^{2}\) , let \(\mathcal{G}_{0}\) be the topology induced by that of \(R^{2}\) , and let \(\mathcal{G}\) be the (finer) topology defined as follows: the neighbourhood filter of a point \(x \neq (1,0)\) of B is the same for \(\mathcal{G}\) as for \(\mathcal{G}_{0}\) , but for the point \((1,0)\) a fundamental system of neighbourhoods in the topology \(\mathcal{G}\) is formed by the sets

\[
\left\{\left(\mathrm{I}, \mathrm{o}\right) \right\} \cup \left(\mathbf {D} \cap \mathbf {B} _ {\mathbf {0}}\right),
\]

where D is any open disc centred at (1,0) and \(B_{0}\) is the open disc \(\|x\|<1\) . Let Y be the set B with the topology C, and let R be the equivalence relation on Y whose equivalence classes are the set S of all \(x\neq(1,0)\)

in B such that \(\|x\| = 1\) and the subsets of Y which consist of a single point of B - S. Show that the quotient space X = Y/R is quasi-compact but not locally quasi-compact. *

\BourbakiExerciseGroup

\begin{enumerate}
\def\labelenumi{\arabic{enumi})}
\tightlist
\item
  \begin{enumerate}
  \def\labelenumii{\alph{enumii})}
  \tightlist
  \item
    Give an example of continuous mappings \(f: \mathbf{X} \to \mathbf{Y}\) , \(g: \mathbf{Y} \to \mathbf{Z}\) such that \(g \circ f\) is proper, \(f\) is not surjective and \(g\) is not proper.
  \end{enumerate}
\end{enumerate}

\begin{enumerate}
\def\labelenumi{\alph{enumi})}
\setcounter{enumi}{1}
\tightlist
\item
  Give an example of two proper mappings \(f: \mathbf{X} \to \mathbf{Y}\) and \(g: \mathbf{X} \to \mathbf{Z}\) where \(\mathbf{X}, \mathbf{Y}, \mathbf{Z}\) are quasi-compact, such that \(x \to (f(x), g(x))\) is not a proper mapping of \(\mathbf{X}\) into \(\mathbf{Y} \times \mathbf{Z}\) .
\end{enumerate}

\begin{enumerate}
\def\labelenumi{\arabic{enumi})}
\setcounter{enumi}{1}
\tightlist
\item
  Let \(f: \mathbf{X} \to \mathbf{Y}\) be a proper mapping and let \(\mathbf{A}\) be a subspace of \(\mathbf{X}\) . Show that if \(\mathbf{A} = \overline{f}^1(f(\mathbf{A})) \cap \overline{\mathbf{A}}\) then the restriction \(f|_{\mathbf{A}}\) is a proper mapping of \(\mathbf{A}\) into \(f(\mathbf{A})\) . Is this condition necessary?
\end{enumerate}

¶ 3) Let \((\mathbf{X}_{\alpha}, f_{\alpha\beta})\) be an inverse system of quasi-compact Kolmogoroff spaces such that the \(f_{\alpha\beta}\) are proper mappings. Show that if the \(\mathbf{X}_{\alpha}\) are non-empty then \(\varprojlim \mathbf{X}_{\alpha}\) is not empty (show that in a quasi-compact Kolmogoroff space, every non-empty closed subset contains a closed subset consisting of a single point).

\begin{enumerate}
\def\labelenumi{\arabic{enumi})}
\setcounter{enumi}{3}
\item
  Give an example of a continuous mapping \(f: \mathbf{X} \to \mathbf{Y}\) (where \(\mathbf{X}\) and \(\mathbf{Y}\) are Hausdorff spaces) such that the inverse image \(\overline{f}^1(\mathbf{K})\) of every compact subset \(\mathbf{K}\) of \(\mathbf{Y}\) is compact, but which is not proper (cf.~§ 9, Exercise 4).
\item
  Let \(f: \mathbf{X} \to \mathbf{Y}\) be a proper mapping.
\end{enumerate}

\begin{enumerate}
\def\labelenumi{\alph{enumi})}
\tightlist
\item
  Show that if X is regular then \(f(\mathbf{X})\) is a regular subspace of Y.
\item
  Show that if \(f(\mathbf{X})\) is a Lindelöf subspace of Y (§ 9, Exercise 14), then X is Lindelöf (use Proposition 10 of § 5, no. 4).
\end{enumerate}

¶ 6) Let X be a space which is the sum of subspaces \(\mathbf{X}_i (i \in I)\) , and let \(f\) be a continuous mapping of X into a topological space Y. For each \(i \in I\) , let \(f_i: \mathbf{X}_i \to \mathbf{Y}\) be the restriction of \(f\) to \(\mathbf{X}_i\) . Show that \(f\) is proper if and only if each of the \(f_i\) is proper and the family \((f(\mathbf{X}_i))_{i \in I}\) is locally finite {[}use Theorem 1 d){]}.

\begin{enumerate}
\def\labelenumi{\arabic{enumi})}
\setcounter{enumi}{6}
\item
  Let \(f: \mathbf{X} \to \mathbf{X}'\) be a proper mapping and let \(\mathbf{R}\) (resp. \(\mathbf{R}'\) ) be an equivalence relation on \(\mathbf{X}\) (resp. \(\mathbf{X}'\) ) with which \(f\) is compatible. Let \(\overline{f}: \mathbf{X} / \mathbf{R} \to \mathbf{X}' / \mathbf{R}'\) be the mapping of the quotient spaces induced by \(f\) . Suppose that (i) the image under \(f\) of every saturated set with respect to \(\mathbf{R}\) is saturated with respect to \(\mathbf{R}'\) , and (ii) \(\mathbf{R}'\) is open. Show that in these conditions \(\overline{f}\) is proper {[}use Theorem 1 c){]}.
\item
  Show that under the hypotheses of Example 3 of § 5, no. 2, the canonical mapping \(\mathbf{X} \to \mathbf{X} / \mathbf{R}\) is proper.
\end{enumerate}

¶ 9) a) Let X, Y be two topological spaces and \(f: \mathbf{X} \to \mathbf{Y}\) a surjective closed (but not necessarily continuous) mapping. Show that if \(\overline{f}^{1}(y)\) is quasi-compact for each \(y \in \mathbf{Y}\) , then \(\overline{f}^{1}(\mathbf{K})\) is quasi-compact for each quasi-compact subset K of Y (either prove the result directly or else use Exercise 7 of § 5 and Exercise 12 of § 9).

\begin{enumerate}
\def\labelenumi{\alph{enumi})}
\setcounter{enumi}{1}
\tightlist
\item
  Suppose in addition that X is regular. Show that if \(\overline{f}^{1}(y)\) is closed in X for each \(y \in Y\) , then \(\overline{f}^{1}(K)\) is closed in X for each quasi-compact subset K of Y (same methods). If also Y is locally compact then f is continuous.
\end{enumerate}

\begin{enumerate}
\def\labelenumi{\arabic{enumi})}
\setcounter{enumi}{9}
\tightlist
\item
  Let X, Y be two topological spaces. A correspondence (G, X, Y) between X and Y (Set Theory, Chapter II, § 3, no. 1) is said to be proper if the restriction of the projection \(\mathbf{pr}_2\) to the graph G is a proper mapping \(\mathbf{G} \to \mathbf{Y}\) .
\end{enumerate}

\begin{enumerate}
\def\labelenumi{\alph{enumi})}
\item
  Let \(f: \mathbf{X} \to \mathbf{Y}\) be a mapping. Show that \(f\) is continuous if and only if the correspondence \(\overline{f}^1\) between \(\mathbf{Y}\) and \(\mathbf{X}\) is proper, and that \(f\) is a proper mapping if and only if the correspondences \(\overline{f}^1\) (between \(\mathbf{Y}\) and \(\mathbf{X}\) ) and \(f\) (between \(\mathbf{X}\) and \(\mathbf{Y}\) ) are proper. Hence construct an example of a non-continuous mapping \(f\) such that the correspondence \(f\) is proper.
\item
  If \((\mathbf{G}, \mathbf{X}, \mathbf{Y})\) is a proper correspondence, show that \(\mathbf{G}(\mathbf{A})\) is a closed subset of Y whenever A is a closed subset of X.
\end{enumerate}

¶ 11) Let (G, X, Y) be a proper correspondence (Exercise 10).

\begin{enumerate}
\def\labelenumi{\alph{enumi})}
\tightlist
\item
  Show that if \(\mathfrak{B}\) is a filter base on \(X\) which consists of closed sets, then
\end{enumerate}

\[
\bigcap_ {\mathbf {M} \in \mathfrak {B}} \mathrm{G} (\mathbf {M}) = \mathrm{G} \left(\bigcap_ {\mathbf {M} \in \mathfrak {B}} \mathbf {M}\right).
\]

\begin{enumerate}
\def\labelenumi{\alph{enumi})}
\setcounter{enumi}{1}
\tightlist
\item
  For each \(y \in \mathbf{Y}\) and each neighbourhood \(\mathbf{V}\) of \(\overline{\mathbf{G}}(y)\) in \(\mathbf{X}\) , show that there is a neighbourhood \(\mathbf{W}\) of \(y\) in \(\mathbf{Y}\) such that \(\overline{\mathbf{G}}^{1}(\mathbf{W}) \subset \mathbf{V}\) (``continuity of the roots as functions of the parameters'') (use Exercise 6 of § 9).
\end{enumerate}

¶ 12) a) A correspondence (G, X, Y) is proper if and only if, for every topological space Z and every correspondence (E, Z, X) whose graph E is closed in Z × X, the graph G \(\circ\) E is closed in Z × Y. First to show that the condition is necessary, remark that G \(\circ\) E is the image of (E × Y) ∩ (Z × G) under \(\iota_{\mathbf{Z}} \times \mathrm{pr}_{2} : \mathbf{Z} \times \mathbf{G} \to \mathbf{Z} \times \mathbf{Y}\) . To show sufficiency, let \(\delta_{\mathbf{Y}}\) be the diagonal mapping \(\mathbf{Y} \to \mathbf{Y} \times \mathbf{Y}\) ; show that, for each subset \(\mathbf{F}\) of \((\mathbf{Z} \times \mathbf{Y}) \times \mathbf{X}\) , the inverse image of \(\mathbf{G} \circ \mathbf{F}\) under the mapping

\[
\iota_ {\mathbf {Z}} \times \delta_ {\mathbf {Y}}: \mathbf {Z} \times \mathbf {Y} \rightarrow \mathbf {Z} \times \mathbf {Y} \times \mathbf {Y}
\]

is the image, under the mapping \((y, z) \to (z, y)\) , of the set

\[
\operatorname{pr} _ {2 3} (\overline {{{\mathrm{F}}}} \cap (\mathrm{G} \times \mathrm{Z})),
\]

where \(\mathbf{F}\) is the image of \(\mathbf{F}\) under the mapping \((z, y, x) \to (x, y, z)\) . b) Deduce from a) that the composition of two proper correspondences is a proper correspondence.

\begin{enumerate}
\def\labelenumi{\alph{enumi})}
\setcounter{enumi}{2}
\tightlist
\item
  Let \((\mathbf{G},\mathbf{X},\mathbf{Y})\) be a proper correspondence. Show that if \(\mathbf{X}\) is Hausdorff then \(\mathbf{G}\) is closed in \(\mathbf{X} \times \mathbf{Y}\) .
\end{enumerate}

¶ 13) Let X be a locally compact space and Y any topological space. a) Show that the following four statements are equivalent:

\(\alpha)\) (G, X, Y) is a proper correspondence between X and Y.

\(\beta)\) For each Hausdorff space \(\mathbf{X}'\) which contains \(\mathbf{X}\) as a subspace, \(\mathbf{G}\) is closed in \(\mathbf{X}' \times \mathbf{Y}\) .

\(\gamma)\) G is closed in \(\mathbf{X} \times \mathbf{Y}\) , and for each \(y \in \mathbf{Y}\) there is a neighbourhood \(\mathbf{V}\) of \(y\) in \(\mathbf{Y}\) and a compact subset \(\mathbf{K}\) of \(\mathbf{X}\) such that \((\mathbf{X} - \mathbf{K}) \times \mathbf{V}\) does not meet \(\mathbf{G}\) .

δ) There is a compact space \(X'\) , containing X as a subspace, such that G is closed in \(X' \times Y\) .

{[}To prove that \(\alpha) \Longrightarrow \beta)\) use Exercise 12 \(a)\) and \(c)\) . To show that \(\delta) \Longrightarrow \alpha)\) , use Proposition 2 of no. 1 and Corollary 5 of Theorem 1 of no. 2.{]}

\begin{enumerate}
\def\labelenumi{\alph{enumi})}
\setcounter{enumi}{1}
\tightlist
\item
  Show that if \((\mathbf{G},\mathbf{X},\mathbf{Y})\) is proper, then:
\end{enumerate}

\(\zeta)\) For every compact subset \(\mathbf{L}\) of \(\mathbf{Y}\) , \(\overline{\mathbf{G}}(\mathbf{L})\) is a compact set.

Show also that if Y is locally compact and if G is closed in \(X \times Y\) and satisfies \(\zeta\) , then (G, X, Y) is proper. {[}Use the equivalence of \(\alpha\) and \(\gamma\) .{]} This last assertion is no longer valid when Y is Hausdorff but not locally compact (Exercise 4).

\begin{enumerate}
\def\labelenumi{\arabic{enumi})}
\setcounter{enumi}{13}
\item
  Let Y be a locally compact space, X a topological space. Show that a mapping \(f: X \rightarrow Y\) is continuous if and only if, for each compact space \(Y'\) which contains Y as a subspace, the graph of f is closed in \(X \times Y'\) . {[}Use Exercises 13 and 10 a){]}.
\item
  Let X be a Hausdorff space, A a closed subset of X, f a proper mapping of X --- A into a locally compact space Y, Y' the one-point compactification of Y, and \(\omega\) the point at infinity in Y'. Let g be the mapping of X into Y' which coincides with f on X --- A and maps each point of A to \(\omega\) . Show that g is continuous.
\end{enumerate}

¶ 16) Let X be a locally compact space, R an equivalence relation on X, C the graph of R in X × X. Show that C is closed in X × X if and only if, for each compact subset K of X, the saturation of K with respect to R is closed in X. (To prove necessity, use Corollary 5 of Theorem 1 of no. 2).

*17) On the locally compact space R, let S be the equivalence relation obtained by identifying all the points of Z. Show that S is closed and R/S Hausdorff but not locally compact.*

*18) Let X be the locally compact subspace {[}o, +∞{[} of R. Let S be the equivalence relation on X whose equivalence classes are the sets \{x, 1/x\} for o \textless{} x ≤ 1 and the set \{o\}. Show that S is open, that the graph of S is closed in X × X and that X/S is compact, but that S is not closed.*

¶ 19) Let X be a locally compact σ-compact space, R an equivalence relation on X whose graph is closed in X × X. Show that X/R is Hausdorff. {[}Let (Uₙ) be a covering of X by relatively compact open sets such that Uₙ ⊂ Uₙ₊₁. Let A, B be two disjoint closed sets which are saturated with respect to R. Define recursively two increasing sequences (Vₙ), (Wₙ) of closed sets saturated with respect to R, such that Vₙ ∩ Wₙ = ∅ and such that Vₙ ∩ Uₙ (resp. Wₙ ∩ Uₙ) is a neighbourhood of A ∩ Uₙ (resp. B ∩ Uₙ) in the compact subspace Uₙ. Use Proposition 8 of no. 3 and Exercise 16, as well as Proposition 2 of § 9, no. 2{]}. (*)

¶ 20) a) Let X be a non-empty compact space and R a Hausdorff equivalence relation on X. Show that X/R is closed in the space \(\mathfrak{P}_{0}(X)\) with the topology \(\mathcal{G}_{\Theta}\) (§ 8, Exercise 12) if and only if the relation R is open. {[}Use Exercise 13 of § 9 and Exercise 16 of § 3 to show that if X/R is closed in \(\mathfrak{P}_{0}(X)\) in the topology \(\mathcal{G}_{\Theta}\) , then the quotient topology on X/R coincides with that induced by \(\mathcal{G}_{\Theta}\) , and then apply Exercise 7 of § 5. Conversely, if R is open, use Exercise 29 of § 8 to show that X/R is closed with respect to \(\mathcal{G}_{\Theta}\) {]}.

*b) Let X be the locally compact subspace of the real plane \(\mathbb{R}^2\) consisting of the points \((\mathtt{I}/n, y)\) , where \(n\) is an integer \(>0\) and \(-\mathtt{I} \leqslant y \leqslant \mathtt{I}\) , and the points \((\mathtt{o}, y)\) where either \(-\mathtt{I} \leqslant y < 0\) or \(0 < y \leqslant \mathtt{I}\) or \(y = 2\) . Let \(p\) be the restriction of \(\mathsf{pr}_1\) to X, and R the equivalence relation associated with \(p\) . Show that R is neither open nor closed, but that X/R is compact and closed in \(\mathfrak{P}_0(X)\) in the topology \(\mathcal{C}_{\Theta}\) . * ¶ 21) (*) a) Let X be an accessible space which is locally quasi-compact (§ 9, Exercise 29). Let \(\mathbf{X}_0\) be the discrete space which has the same underlying set as X, and let \(\tilde{\mathbf{X}}_0\) be the (compact) ultrafilter space of \(\mathbf{X}_0\) (the discrete space \(\mathbf{X}_0\) being identified with an open subspace of \(\tilde{\mathbf{X}}_0\) ) (§ 9, Exercise 26). On \(\tilde{\mathbf{X}}_0\) , let R {[}or R(X){]} be the following equivalence relation between ultrafilters U, B: ``the sets of limit points of U and B in X are the same''. Show that R is Hausdorff {[}use the description of the topology of \(\tilde{\mathbf{X}}_0\) given in § 9, Exercise 25 b), Exercise 7 b) of § 7 and axiom (C'){]}. The compact space \(\mathbf{X}^c = \tilde{\mathbf{X}}_0 / \mathbf{R}\) is in one-to-one correspondence with the set of primitive subsets of X (§ 7, Exercise 7). If X is Hausdorff (and therefore locally compact) then \(\mathbf{X}^c\) can be identified with X if X is compact, and with the one-point compactification of X if X is not compact.

\begin{enumerate}
\def\labelenumi{\alph{enumi})}
\setcounter{enumi}{1}
\item
  Let Y be a closed subspace of X. The ultrafilter space \(\tilde{Y}_{0}\) can be identified with the closure of \(Y_{0}\) in \(\tilde{X}_{0}\) , and \(Y^{c}\) with the canonical image of \(\tilde{Y}_{0}\) in \(X^{c}\) .
\item
  Show that the restriction to \(X_{0}\) of the canonical mapping \(\tilde{X}_{0} \rightarrow X^{c}\) is a bijection of \(X_{0}\) onto a subspace \(X_{*}^{c}\) of \(X^{c}\) , that the canonical mapping \(f: X_{*}^{c} \rightarrow X\) is a continuous bijection, and that the image under this mapping of any compact subset of \(X_{*}^{c}\) is a closed compact subset of X. {[}To show that f is continuous, use b); also remark that the inverse image in \(\tilde{X}_{0}\) of a point of \(X_{*}^{c}\) consists of all the ultrafilters on X which have a single given limit point, and that if U is an ultrafilter on X, then its limit point in \(\tilde{X}_{0}\) , when U is regarded as an ultrafilter base on \(\tilde{X}_{0}\) , is U regarded as a point of \(\tilde{X}_{0}\) {]}.
\end{enumerate}

*d) Let \(\mathbf{Z}\) be the set which is the sum of a countable number of spaces \(\mathbf{Y}_n\) each identical with the space \(\mathbf{Y}\) defined in Exercise 7 of §8 and a set \(\{\omega\}\) consisting of a single point. Define a topology on \(\mathbf{Z}\) by taking as a fundamental system of neighbourhoods of a point \(z \in \mathbf{Y}_n\) its neighbourhood filter in the space \(\mathbf{Y}_n\) , and as fundamental system of neighbourhoods of \(\omega\) the set of subsets \(\mathbf{V}_n\) of \(\mathbf{Y}\) , where \(\mathbf{V}_n\) is the union of \(\{\omega\}\) and the \(\mathbf{Y}_m\) where \(m \geqslant n\) . The space \(\mathbf{Z}\) so defined is accessible, quasi-compact and locally quasi-compact; but the subspace \(Z_*^c\) of \(Z^c\) is not locally compact.*

\begin{enumerate}
\def\labelenumi{\alph{enumi})}
\setcounter{enumi}{4}
\tightlist
\item
  Let Y be another locally quasi-compact accessible space, and let u be a continuous mapping of X into Y, such that the inverse image under u of each quasi-compact subset of Y is quasi-compact. Regarding u as a mapping of \(X_{*}^{c}\) into \(Y_{*}^{c}\) , show that u is continuous and can be (uniquely) extended to a continuous mapping \(\mathbf{X}^c\to \mathbf{Y}^c\) . {[}First extend \(u\) to a continuous mapping of \(\tilde{\mathbf{X}}_0\) into \(\tilde{\mathbf{Y}}_0\) and show that the extension is compatible with the equivalence relations \(\mathbf{R}(\mathbf{X}),\mathbf{R}(\mathbf{Y})\) {]}.
\end{enumerate}

\(^{*}f)\) Let X be the quasi-compact, not locally quasi-compact, accessible space defined in Exercise 29 c) of § 9. Show that in the corresponding space \(\tilde{X}_{0}\) the equivalence relation R(X) is not Hausdorff.*

\BourbakiExerciseGroup

\begin{enumerate}
\def\labelenumi{\arabic{enumi})}
\item
  Show that the countable Hausdorff space defined in § 9, Exercise 20 c), is connected.
\item
  \begin{enumerate}
  \def\labelenumii{\alph{enumii})}
  \tightlist
  \item
    Let X be a topological space and \(\mathfrak{C}\) the topology of X. Show that if X is connected in the semi-regular topology \(\mathfrak{C}^*\) associated with \(\mathfrak{C}\) ( \(\S 8\) , Exercise 20), then X is connected in the original topology \(\mathfrak{C}\) . Deduce that there exist connected submaximal spaces ( \(\S 8\) , Exercise 22).
  \end{enumerate}
\end{enumerate}

*b) In particular construct a topology on the real line \(\mathbf{R}\) , with respect to which \(\mathbf{R}\) is connected but no point of \(\mathbf{R}\) has a fundamental system of connected neighbourhoods {[}cf.~Chapter IV, § 2, Exercise 14 c){]}.*

\begin{enumerate}
\def\labelenumi{\arabic{enumi})}
\setcounter{enumi}{2}
\tightlist
\item
  Let A, B be two subsets of a topological space X.
\end{enumerate}

\begin{enumerate}
\def\labelenumi{\alph{enumi})}
\item
  Show that if A and B are closed in X, and A ∪ B and A ∩ B are connected, then A and B are connected. Show by an example that this statement is no longer necessarily true when one of the sets A, B is not closed.
\item
  Suppose that A and B are connected and \(\overline{A} \cap B \neq \varnothing\) . Show that \(A \cup B\) is connected.
\end{enumerate}

\begin{enumerate}
\def\labelenumi{\arabic{enumi})}
\setcounter{enumi}{3}
\tightlist
\item
  Let X be a connected space of at least two points.
\end{enumerate}

\begin{enumerate}
\def\labelenumi{\alph{enumi})}
\item
  Let A be a connected subset of X, B a subset of \(\mathbb{C}\mathbf{A}\) which is both open and closed in \(\mathbb{C}\mathbf{A}\) . Show that \(\mathbf{A} \cup \mathbf{B}\) is connected. (Apply Exercise 5 of §3 to \(\mathbf{Y} = \mathbf{A} \cup \mathbf{B}\) and \(\mathbf{Z} = \mathbb{C}\mathbf{A}\) .)
\item
  Let A be a connected subset of X and B a component of the set CA. Show that CB is connected {[}use a){]}.
\item
  Deduce from b) that there exist two connected subsets M, N of X which are distinct from X and such that \(M \cup N = X\) and \(M \cap N = \varnothing\) .
\end{enumerate}

*5) a) Give an example, in the real plane \(\mathbf{R}^2\) , of a decreasing sequence \((\mathbf{A}_n)\) of connected sets whose intersection is not connected (cf.~Chapter II, § 4, Exercise 14).

\begin{enumerate}
\def\labelenumi{\alph{enumi})}
\setcounter{enumi}{1}
\tightlist
\item
  In \(\mathbf{R}^2\) , let \(\mathbf{X}\) be the set which is the union of the open half-planes \(y > 0\) and \(y < 0\) and the points \((n, 0)\) where \(n\) is an integer \(> 0\) .
\end{enumerate}

X is connected with respect to the topology \(\mathfrak{C}_0\) induced by that of \(\mathbb{R}^2\) . Define a sequence \((\mathfrak{C}_n)\) of topologies on X such that \(\mathfrak{C}_{n+1}\) is finer than \(\mathfrak{C}_n\) , and such that X is connected with respect to each \(\mathfrak{C}_n\) but not connected with respect to the topology which is the least upper bound of the sequence \((\mathfrak{C}_n)\) {[}pass from \(\mathfrak{C}_n\) to \(\mathfrak{C}_{n+1}\) by introducing new neighbourhoods of the point \((n+1,0)]\) .*

\begin{enumerate}
\def\labelenumi{\arabic{enumi})}
\setcounter{enumi}{5}
\item
  Let X, Y be two connected spaces and let A (resp. B) be a subset of X (resp. Y) distinct from X (resp. Y). Show that in the product space \(X \times Y\) the complement of \(A \times B\) is connected.
\item
  Let X, Y be two connected spaces and f a mapping of \(X \times Y\) into a topological space Z such that all the partial mappings \(f(0, y) : X \to Z\) , \(f(x, 0) : Y \to Z (x \in Y, y \in Y)\) are continuous. Show that \(f(X \times Y)\) is connected.
\item
  In the example of a topology on a product set \(\mathbf{X} = \prod_{i\in I}\mathbf{X}_i\) given in § 4, Exercise 9, suppose that I is infinite and the cardinal t uncountable. Show that if the \(\mathbf{X}_i\) are connected, regular, and have each at least two distinct points, then the connected component of a point \(a = (a_i)\) of X is the set of \(x = (x_i)\) such that \(x_{i} = a_{i}\) except for a finite number of indices \(i\in I\) . (Reduce to the case I = N; consider two points \(b = (b_n)\) , \(c = (c_n)\) of X such that \(b_{n}\neq c_{n}\) for all n.~For each n let \((\mathrm{V}_{nm})_{m\geqslant 0}\) be a sequence of open neighbourhoods of \(b_{n}\) in \(\mathbf{X}_n\) which do not contain \(c_{n}\) and are such that \(\overline{\mathbf{V}}_{n,m + 1}\subset \mathbf{V}_{nm}\) . Let Z be the set of points \((x_{n})\) of X with the following property: there exists an integer \(k > 0\) such that \(x_{n}\in \mathrm{V}_{n,n - k}\) whenever \(n\geqslant k\) . Show that Z is both open and closed in X).
\end{enumerate}

*9) Show that in the locally compact space \(\mathbf{X}\) defined in § 10, Exercise 20, there are points \(x\) whose component in \(\mathbf{X}\) is different from the intersection of the subsets of \(\mathbf{X}\) which are both open and closed and contain \(x_*\)

\begin{enumerate}
\def\labelenumi{\arabic{enumi})}
\setcounter{enumi}{9}
\item
  Show that a linearly ordered set X, with the topology \(\mathcal{G}_{+}(X)\) or the topology \(\mathcal{G}_{-}(X)\) (§ 2, Exercise 5) is totally disconnected.
\item
  Let X be a topological space and R an equivalence relation on X such that every equivalence class mod R is contained in a component of X. Show that the components of X/R are the canonical images of the components of X (cf.~Chapter III, § 2, Exercise 17).
\item
  Let X be a topological space. Show that the following three conditions are equivalent: \(\alpha\) ) the components of X are open sets; \(\beta\) ) every \(x \in X\) has a neighbourhood contained in every set which is both open and closed and contains \(x; \gamma\) ) for every \(x \in X\) , the intersection of the sets which are both open and closed and contain x is an open set.
\end{enumerate}

*13) Let \(\theta\) be an irrational number. Let \(f\) be the mapping of the real line \(\mathbf{R}\) into the torus \(\mathbf{T}^2\) for which the image \(f(t)\) of each \(t \in \mathbf{R}\) is the canonical image in \(\mathbf{T}^2\) of \((t, \theta t) \in \mathbf{R}^2\) . \(f\) is injective and continuous. Show that, in the topology on \(\mathbf{R}\) which is the inverse image under \(f\) of the topology of \(\mathbf{T}^2\) , \(\mathbf{R}\) is connected but no point of \(\mathbf{R}\) has a fundamental system of connected neighbourhoods.*

\begin{enumerate}
\def\labelenumi{\arabic{enumi})}
\setcounter{enumi}{13}
\tightlist
\item
  \begin{enumerate}
  \def\labelenumii{\alph{enumii})}
  \tightlist
  \item
    Let X be a locally compact and locally connected space, and let A be a closed subset of X. Let A* be the union of A and the components of CA which are relatively compact. A is said to be full if A* = A. Show that A* is closed and full.
  \end{enumerate}
\end{enumerate}

\begin{enumerate}
\def\labelenumi{\alph{enumi})}
\setcounter{enumi}{1}
\item
  Suppose in addition that X is connected. Show that if A is a compact subset of X, then A* is compact (consider a compact neighbourhood V of A and the components of CA which meet the frontier of V).
\item
  Suppose in addition that X is connected and \(\sigma\) -compact. Show that there exists an increasing sequence \((\mathbf{K}_{n})_{n\geqslant0}\) of full connected compact sets whose union is X and which are such that \(K_{n}\subset\mathring{K}_{n+1}\) for all n.
\end{enumerate}

*d) Show that the conclusion of b) is no longer valid if we suppose only that X is connected and locally compact, or that X is Hausdorff, connected and locally connected. {[}Consider the subspace of \(\mathbf{R}^2\) consisting of points \((x, y)\) such that either \(y = 0\) and \(0 \leqslant x \leqslant 1\) , or else \(y > 0\) and \(x = 1/n\) where \(n\) is an integer \(> 0\) .*

¶ 15) Let X be a connected, locally connected space.

\begin{enumerate}
\def\labelenumi{\alph{enumi})}
\item
  Let M, N be two disjoint non-empty closed subsets of X. Show that \(\mathbb{C}(M \cup N)\) has a component whose closure meets both M and N. {[}If C is a component of \(\mathbb{C}(M \cup N)\) whose closure does not meet M, observe that C is both open and closed in CN{]}.
\item
  If A is a closed subset of X distinct from X, deduce from a) that every component of A meets the closure of \(\mathbb{C}\mathbb{A}\) .
\end{enumerate}

*16) Let X be the subspace of the real plane \(\mathbf{R}^2\) which is the union of the line \(y = 0\) , the closed segments with end points (0, 1) and \((n, 1/(1 + n))\) , and the half-lines \(x \leqslant n\) , \(y = 1/(1 + n)\) ( \(n\) an integer \(>0\) ). Show that X is a connected space but that there is a component C of the complement of (0, 1) in X such that (0, 1) \(\notin \overline{\mathbf{C}}\) {[}(compare with Exercises 4 a) and 15 b){]}.*

\(^{*}\) 17) Let \(I_{0}\) be the interval \([-1,1]\) of R and let \(P_{0}\) be the quotient space of \(I_{0}\) obtained by identifying the two points 1/2 and 1; let \(B_{0}\) be the quotient space of \(I_{0}\) obtained by identifying the points 1/2 and 1 and also identifying -1/2 and -1. Let I be the open interval {]}--- I, I{[} in R, and let P (resp. B) be the canonical image of I in \(P_{0}\) (resp. \(B_{0}\) ).

\begin{enumerate}
\def\labelenumi{\alph{enumi})}
\item
  Show that no two of the three spaces I, P, B are homeomorphic (to show that I and P are not homeomorphic, notice that there are points in P whose complement is connected; similarly for the other pairs of spaces).
\item
  Let X be the sum of a countably infinite family of spaces homeomorphic to I and a countably infinite family of spaces homeomorphic to B. Let Y be the sum of X and P. Show that X and Y are not homeomorphic, but that (i) there is a continuous bijection of X onto Y and a continuous bijection of Y onto X; (ii) X is homeomorphic to an open subspace of Y and Y is homeomorphic to an open subspace of X.*
\end{enumerate}

*18) Let X be the real plane \(\mathbf{R}^2\) , and let R be the equivalence relation on X whose equivalence classes are the lines \(y = \beta\) for \(\beta > 0\) and the half-lines \(x = \alpha\) , \(y \leqslant 0\) for all \(\alpha \in \mathbb{R}\) . Show that if X/R is regarded as a subset of \(\mathfrak{P}_0(X)\) , then the topology induced on X/R by \(\mathcal{C}_{\Omega}\) ( \(\S 2\) , Exercise 7) is discrete, and that the topology induced on X/R by \(\mathcal{C}_{\Phi}\) is not discrete but makes X/R into a non-connected space. Show finally that X/R is connected in the quotient topology.*

*¶ 19) Let X be a locally compact space. It can be shown that X can be identified with a dense open subspace of the universal compact space \(\tilde{X}\) constructed in § 9, Exercise 25 (cf.~Chapter IX, § 1, Exercise 8). Suppose also that X is connected and locally connected.

\begin{enumerate}
\def\labelenumi{\alph{enumi})}
\item
  Show that, for every compact subset K of X, the set of components of X-K which are not relatively compact is finite {[}cf.~Exercise 14 b){]}. Deduce that a point of \(\tilde{X} - X\) belongs to the closure of at most one of the components of X-K {[}if \(A_i\) ( \(i \leqslant i \leqslant n\) ) are the non-compact components of X-K, and \(C_i\) the compact set of points of \(\tilde{X} - X\) which lie in the closure of \(A_i\) , construct a compact space whose underlying set is the sum of X and the sets \(C_i\) , such that X is identified with a dense open subspace of this space{]}.
\item
  Let \(R'\{x, y\}\) be the following equivalence relation on \(\tilde{X}-X\) : ``for every compact subset K of X, x and y lie in the closure of the same non-compact component of X-K''. Let R be the equivalence relation on \(\tilde{X}\) whose equivalence classes are the classes mod \(R'\) and the subsets consisting of a single point of X. Show that R is Hausdorff and that X may be identified with a dense open subspace of the compact space \(X^{b}=\tilde{X}/R\) ; show that in the compact space \(X^{b}-X\) every point is the intersection of its neighbourhoods which are both open and closed. The points of \(X^{b}-X\) are called the ends of the space X.
\item
  Let Y be a compact space such that X may be identified with a dense open subspace of Y and such that in the compact space Y-X every point is the intersection of its neighbourhoods which are both open and closed (*). Show that there is a continuous mapping of \(X^{b}\) into Y which induces the identity mapping on X. {[}If Y-X is the union of two disjoint non-empty closed sets A and B, show that there exist two disjoint open subsets U, Y in V such that \(A \subset U\) , \(B \subset V\) and such that \(Y - (U \cup V)\) is compact{]}.
\item
  Show that if X is σ-compact, then \(X^{b}-X\) has a countable base.
\item
  The real line has two ends, and \(R^{n} (n \geqslant 2)\) has one. Give an example of a connected, locally connected, locally compact subspace of \(R^{2}\) for which the set of ends has the power of the continuum.*
\end{enumerate}

\begin{enumerate}
\def\labelenumi{\arabic{enumi})}
\setcounter{enumi}{19}
\tightlist
\item
  Let X be a Hausdorff space in which each point has a fundamental system of neighbourhoods which are both open and closed in X. Let Y be a subspace of X, and let A be a compact subset of Y which is both open and closed in Y; show that there is a subset B of X which is both open and closed in X, such that \(B \cap Y = A\) .
\end{enumerate}

¶ 21) a) Show that the following conditions on a topological space X are equivalent: (i) for each open set U in X, \(\overline{U}\) is open; (ii) for each closed set F in X, \(\dot{F}\) is closed; (iii) for each disjoint pair of open sets U, V in X, we have \(\overline{U} \cap \overline{V} = \varnothing\) . A Hausdorff space which satisfies these conditions is said to be extremely disconnected; it is then totally disconnected. *The rational line is totally disconnected but not extremely disconnected.*

\begin{enumerate}
\def\labelenumi{\alph{enumi})}
\setcounter{enumi}{1}
\item
  A Hausdorff space is extremely disconnected if and only if the associated semi-regular space (§ 8, Exercise 20) is extremely disconnected.
\item
  Let X be a Hausdorff space and let A be a dense subset of X. Show that if A is extremely disconnected and if the topology of X is A-maximal (§ 3, Exercise 11), then X is extremely disconnected.
\item
  Deduce from b) and c) that, if X is any infinite set, the (compact) ultrafilter space of X is extremely disconnected (cf.~§ 9, Exercises 26 and 24) (**).
\item
  Let X be a Hausdorff space with no isolated points. Show that the topology of X is quasi-maximal (§ 2, Exercise 6) if and only if X is submaximal (§ 8, Exercise 22) and extremely disconnected. (To show that the condition is sufficient, use Exercise 22 e) of § 8 and prove that every closed subset of X which has no isolated points is open).
\item
  Show that every ultraregular space (§ 8, Exercise 25) is extremely disconnected, but that an extremely disconnected compact space with no isolated points is not ultraregular (cf.~§ 9, Exercise 27).
\item
  Show that an extremely disconnected semi-regular space (§ 8, Exercise 20) is regular (cf.~Chapter II, § 4, Exercise 12).
\item
  Show that in an extremely disconnected space there is no convergent sequence which has an infinite number of distinct terms {[}argue by contradiction by constructing recursively two sequences \((\mathbf{U}_{n})\) , \((\mathbf{V}_{n})\) of mutually disjoint open sets such that if \(U = \bigcup_{n} U_{n}\) and \(V = \bigcup_{n} V_{n}\) then \(U \cap V = \varnothing\) but \(\overline{U} \cap \overline{V} \neq \varnothing\) {]}.
\item
  Show that in an extremely disconnected space X a non-isolated point a cannot have a fundamental system of simultaneously open and closed neighbourhoods which is well-ordered (with respect to the relation \(\supset\) ). (Use reductio ad absurdum, as in h).
\end{enumerate}

¶ 22) a) Show that in an extremely disconnected space every open subspace and every dense subspace is extremely disconnected.

\begin{enumerate}
\def\labelenumi{\alph{enumi})}
\setcounter{enumi}{1}
\item
  Let X be a Hausdorff space in which every point has a neighbourhood which is an extremely disconnected subspace of X. Show that X is extremely disconnected.
\item
  Give an example of a Hausdorff space X which is not extremely disconnected but which has an extremely disconnected dense open subspace (paste together two extremely disconnected spaces).
\item
  Let X be a countable discrete space, \(\tilde{X}\) the (compact) space of ultrafilters on X (§ 9, Exercise 26). Let \((\mathbf{A}_{n})_{n\in\mathbb{N}}\) be a countably infinite partition of X into infinite subsets; in the closed subspace \(Y=\tilde{X}-X\) of \(\tilde{X}\) , the sets \(B_{n}=\tilde{X}_{n}\cap Y\) are open and closed and pairwise disjoint. Let \(B=\bigcup_{n}B_{n}\) ; show that the closure B of B in Y is not open in Y, and hence that Y is not extremely disconnected (*). (Use reductio ad absurdum: using Exercise 20, let C be an open and closed subset of \(\tilde{X}\) such that \(C\cap Y=B\) ; consider a point \(x_{n}\) in each of the sets \(C\cap A_{n}\) , let \(\bar{J}\) be the closure of the set J of points \(x_{n}\) , and show that \(\bar{J}\) does not meet B).
\end{enumerate}

\begin{enumerate}
\def\labelenumi{\arabic{enumi})}
\setcounter{enumi}{22}
\item
  Give an example of a bijective mapping \(f\) of a Hausdorff space \(\mathbf{X}\) onto a Hausdorff space \(\mathbf{Y}\) , such that the image under \(f\) of every compact subset of \(\mathbf{X}\) is a compact subset of \(\mathbf{Y}\) , and the image under \(f\) of every connected subset of \(\mathbf{X}\) is a connected subset of \(\mathbf{Y}\) , but such that \(f\) is not continuous (cf.~§ 9, Exercise 4).
\item
  Let X be a topological space.
\end{enumerate}

\begin{enumerate}
\def\labelenumi{\alph{enumi})}
\tightlist
\item
  Let the set \(\mathfrak{P}_{\Omega}(\mathbf{X})\) carry one of the topologies \(\mathfrak{C}_{\Omega}, \mathfrak{C}_{\Phi}\) ( \(\S 2\) , Exercise 7) or \(\mathfrak{C}_{\Theta}\) ( \(\S 8\) , Exercise 12). Show that if a subset \(\mathfrak{B}\) of \(\mathfrak{P}_{\Omega}(\mathbf{X})\) is connec-
\end{enumerate}

ted and if each \(\mathbf{M} \in \mathfrak{B}\) is connected, then \(\bigcup_{\mathbf{M} \in \mathfrak{M}} \mathbf{M}\) is connected.

\begin{enumerate}
\def\labelenumi{\alph{enumi})}
\setcounter{enumi}{1}
\tightlist
\item
  Let \(\mathfrak{S}\) be a subset of \(\mathfrak{P}_{0}(\mathbf{X})\) which contains the set of all finite non-empty subsets of \(\mathbf{X}\) . Show that if \(\mathbf{X}\) is connected then so is \(\mathfrak{S}\) {[}use Exercise 7 b) of § 2 and Proposition 8 of § 4, considering the mappings \((x_{1},\ldots ,x_{n})\to \{x_{1},\ldots ,x_{n}\} ]\) .
\end{enumerate}

\begin{enumerate}
\def\labelenumi{\arabic{enumi})}
\setcounter{enumi}{24}
\tightlist
\item
  Given two topological spaces X, Y, a mapping \(f: \mathbf{X} \to \mathbf{Y}\) is said to be a local homeomorphism if for each \(x \in \mathbf{X}\) there is a neighbourhood U of \(x\) such that \(f(\mathbf{U})\) is a neighbourhood of \(f(x)\) in Y, and the mapping \(\mathbf{U} \to f(\mathbf{U})\) which coincides with \(f\) on U is a homeomorphism. Every local homeomorphism is a continuous open mapping.
\end{enumerate}

\begin{enumerate}
\def\labelenumi{\alph{enumi})}
\item
  Suppose that X is Hausdorff and Y locally connected. Let \(f: \mathbf{X} \to \mathbf{Y}\) be a local homeomorphism with the property that there is an integer \(n > 0\) such that \(\overline{f}^1(y)\) has exactly \(n\) points for every \(y \in \mathbf{Y}\) . Show that every \(y \in \mathbf{Y}\) has an open neighbourhood V such that \(\overline{f}^1(V)\) has \(n\) components \(\mathbf{U}_i\) ( \(1 \leqslant i \leqslant n\) ) and such that the mapping \(\mathbf{U}_i \to \mathbf{V}\) which coincides with \(f\) on \(\mathbf{U}_i\) is a homeomorphism. \(f\) is then a proper mapping.
\item
  Suppose that X is Hausdorff and that Y is connected and locally connected. Let \(f: \mathbf{X} \to \mathbf{Y}\) be a proper local homeomorphism. Show that \(f\) satisfies the hypotheses of \(a\) ). {[}If for each integer \(n > 0\) , \(\mathbf{Y}_n\) is the set of \(y \in \mathbf{Y}\) such that \(\bar{f}^1(y)\) has at least \(n\) points, show that \(\mathbf{Y}_n\) is both open and closed in \(\mathbf{Y}\) {]}.
\end{enumerate}

\(^{*}c)\) Give an example of a surjective local homeomorphism \(f: X \to Y\) , in which Y is compact, connected and locally connected, X is locally compact, connected and locally connected, \(\overline{f}^{1}(y)\) consists of at most two points for each \(y \in Y\) , but f is not proper (cf.~§ 5, Exercise 4).
